% !TeX root = ../../Book2.tex
% 纲要标题：## 第16章　Morita 理论
% 纲要位置：计划纲要.md，第 1562 行。

\chapter{Morita 理论}
\label{b2:ch:16}
\BookChapterNumber{b2:ch:16}

\begin{chapterabstract}
本章证明模范畴的等价如何由有限生成投射生成元实现，并核对张量函子、Hom 函子与自同态反环的方向。矩阵环和满幂等角环给出具体等价，进一步辨认等价所保持的模论性质及环本身可能改变的性质。
\end{chapterabstract}

\begin{learninggoals}
\begin{itemize}
\item 区分作为对象的生成元、投射性与有限生成性，证明生成元的三种等价刻画。
\item 在左模约定下正确使用 \(S=\operatorname{End}_R(P)^{\mathrm{op}}\)，写出伴随单位、余单位、互逆双模和有限对偶基公式。
\item 从给定范畴等价恢复有限生成投射生成元，并用任意自由展示证明反向的 Morita 等价。
\item 计算矩阵环与满角环的等价，辨认所保留的模结构及可以改变的坐标、维数和生成元数。
\end{itemize}
\end{learninggoals}

设 \(k\) 为域、\(n\ge1\)，令 \(A=M_n(k)\)，并在一个左 \(A\) 模 \(M\) 上取矩阵单位 \(E_{ij}\)。元素 \(m\) 可拆成 \(\sum_iE_{ii}m\)，而 \(E_{ij}\) 把第 \(j\) 个分量送往第 \(i\) 个分量。知道 \(E_{11}M\) 后，其余分量似乎都能由这些矩阵单位重新造出；模同态也由它在这一分量上的作用决定。这说明环的元素和乘法不同，并不妨碍它们的全部模具有相同的组织方式。

\ProseParagraphBreak

从一个矩阵角回收整个模，需要说明重建映射在任意模上都互逆，而且对模同态相容。一般环未必有现成的矩阵单位，这时要问是否存在一个模 \(P\)，从它出发仍能恢复整个模范畴。生成元性质要求每个模都是 \(P\) 某个直和的商；投射性使 \(\operatorname{Hom}_R(P,-)\) 保持正合列；有限生成性使这个函子保持任意直和。三者各自保护一项重建所需的操作，由此构造的 Hom 与张量函子才可能在全部模上互为拟逆。

\ProseParagraphBreak

本章需要第五章的范畴等价、第六章的双模张量与 Hom–张量伴随，以及第七章的投射模和有限对偶基。第十一章的角环自同态公式与第十五章的基本角代数提供具体应用；中心和有限长度的讨论还使用前两编已建立的自然性与合成列。等价始终指全部左模范畴的等价，因而还须构造在一般模上的作用，有限自由模上的检验并不充分。

\ProseParagraphBreak
一般环的结构定理及矩阵环例子可参阅 Weibel 的《The K-book》\cite[Chapter II, Theorem 2.7, Example 2.7.2]{Weibel2013KBook}。该书采用右模，本章固定左模并相应使用自同态反环。有限维代数的版本可参阅 Etingof 等人的《Introduction to Representation Theory》\cite[Theorem 9.6.4, Problem 9.6.5, p. 215; Section 9.7, p. 216]{EtingofEtAl2011RepresentationTheory}；那里的有限维模范畴与本章全部模范畴的范围须加以区分。

% !TeX root = ../../Book2.tex
% 纲要标题：### 16.1 模范畴与生成元
% 纲要位置：计划纲要.md，第 1564 行。

\section{模范畴与生成元}
\label{b2:sec:1601}

上一编把矩阵环 \(M_n(D)\) 的模分解为简单列模的直和。当 \(n>1\) 时，一个列模比正则模小，仍可用它的直和与商构造其他模。对于一般环，也要问哪些模能够这样替代正则模，产生整个模范畴。本节把这种性质与投射性、有限生成性分开说明。始终使用左模；\(R\text{-}\mathbf{Mod}\) 包含全部幺性左 \(R\) 模，而不只包含有限生成模。后面的等价也将在全部模上成立。

% 纲要要点（供目录核对）：
%
% * 投射生成元；区分模忠实、非零对象检测与同态检测，正式证明迹理想判据及对象检测反例
% * 有限生成投射生成元；用两组有限直和因子分别控制有限投射与生成，正式给出三条件各自缺失的反例
% * 矩阵环与原环的模范畴；矩阵单位读取与恢复分量，正则模未必与正则模对应
%

\subsection{投射生成元}
\label{b2:subsec:1601:01}

正则模通过自由模的满射产生所有模。若想用一个模 \(P\) 替代它，最直接的要求就是每个模都能由若干份 \(P\) 的直和取商得到。“生成元”的名称表达了这个要求；下面用 Hom 函子的忠实性定义它，并证明这两种说法等价。

\begin{definition}\label{b2:def:module-generator}
设 \(R\) 为含幺环，\(P\) 为幺左 \(R\) 模。若函子 \(\operatorname{Hom}_R(P,-)\) 忠实，即对任意幺左 \(R\) 模 \(M,N\) 及不同的 \(R\) 线性映射 \(f,g:M\rightarrow N\)，都存在 \(R\) 线性映射 \(h:P\rightarrow M\) 使 \(fh\ne gh\)，则称 \(P\) 为左 \(R\) 模范畴的\term[mo-fanchou-shengchengyuan]{生成元}[generator]。若 \(P\) 还是投射模，则称它为\term[toushe-shengchengyuan]{投射生成元}[projective generator]。
\end{definition}

这里的“生成元”是一个模对象，不是某个模内部的一个元素。条件 \(M\ne0\Rightarrow\operatorname{Hom}_R(P,M)\ne0\) 只检测非零对象，定义中的忠实性则比较不同同态；二者的差别见\cref{b2:prop:object-detection-not-generator}的双数例子。正则模 \(R\) 是生成元：若 \(f(m)\ne g(m)\)，取 \(h(r)=rm\)，在一处求值便可检测两同态不同。

\ProseParagraphBreak
每个模 \(M\) 本来就有自由满射 \(R^{(I)}\to M\)，所以只要 \(P\) 的直和能满射到 \(R\)，便能对自由模的各个分量作同样构造，再满射到 \(M\)。而 \(R\) 由一生成，一的原像又会把这个构造缩到有限份 \(P\)；这就是下面第三个条件的来源。

\begin{theorem}[生成元的等价刻画]\label{b2:thm:generator-characterizations}
设 \(R\) 为含幺环，\(P\) 为幺左 \(R\) 模，所论模均为幺左 \(R\) 模。以下条件等价：
\begin{equivalences}
\item\label{b2:item:generator-faithful} \(\operatorname{Hom}_R(P,-)\) 忠实。
\item\label{b2:item:generator-quotients} 对每个左 \(R\) 模 \(M\)，都存在指标集 \(I\) 及满同态 \(P^{(I)}\rightarrow M\)。
\item\label{b2:item:generator-finite-summand} 存在有限整数 \(n\geq0\)，使正则左模 \({}_RR\) 是 \(P^n\) 的直和因子。
\end{equivalences}
这些等价不要求 \(P\) 已经投射或有限生成。
\end{theorem}
\begin{proof}
先设 Hom 函子忠实。对任意 \(M\)，将全部同态 \(h:P\rightarrow M\) 作为指标，定义求值映射
\[
q:\bigoplus_{h\in\operatorname{Hom}_R(P,M)}P\longrightarrow M,
\qquad (p_h)_h\longmapsto\sum_h h(p_h).
\]
若像 \(T\) 不等于 \(M\)，商映射 \(\pi:M\rightarrow M/T\) 非零，却满足 \(\pi h=0\) 对每个 \(h\) 成立。于是 \(\operatorname{Hom}_R(P,\pi)\) 与零同态诱导的映射相同，违反忠实性。因此 \(q\) 满射，得到第二项。

将第二项用于 \(M=R\)。一的某个原像只有有限个非零坐标，故限制到相应的 \(P^n\) 后，映射 \(q_0:P^n\rightarrow R\) 的像仍包含一，因而包含每个 \(r=r1\)，所以仍满射。这里有限的 \(n\) 来自直和元素的有限支集，不要求 \(P\) 有限生成。取 \(v\in P^n\) 使 \(q_0(v)=1\)，则 \(s(r)=rv\) 为左线性截面，\(q_0s=\operatorname{id}_R\)。所以 \(R\) 是 \(P^n\) 的直和因子。

若第三项成立，任意模 \(M\) 的自由满射 \(R^{(I)}\rightarrow M\) 可以与各分量上的分裂满射 \(P^n\rightarrow R\) 复合，得到某个 \(P\) 直和到 \(M\) 的满射，所以第二项成立。最后，若 \(f\ne g:M\rightarrow N\)，而每个 \(h:P\rightarrow M\) 都满足 \(fh=gh\)，则对第二项给出的满射 \(P^{(I)}\rightarrow M\)，两复合在每个直和分量上相同，因而整体相同。满射性迫使 \(f=g\)，矛盾。故 Hom 函子忠实。
\end{proof}

生成元必为忠实模，即 \(\operatorname{Ann}_R(P)=0\)：若 \(rP=0\)，第三项使 \(r\) 也零化直和因子 \(R\)，故 \(r=r1=0\)。但模的零化子为零，弱于 Hom 函子忠实。例如 \(\mathbb Q\) 作为 \(\mathbb Z\) 模忠实，但 \(\operatorname{Hom}_{\mathbb Z}(\mathbb Q,\mathbb Z)=0\)：任意同态的每个值都被所有正整数整除，只能为零。因此任何 \(\mathbb Q\) 的直和都不能满射到 \(\mathbb Z\)，它不是生成元。

\subsection{有限生成投射生成元}
\label{b2:subsec:1601:02}

\begin{definition}\label{b2:def:progenerator}
设 \(R\) 为含幺环，\(P\) 为幺左 \(R\) 模。若 \(P\) 有限生成、投射，且是左 \(R\) 模范畴的生成元，则称 \(P\) 为\term[youxian-toushe-shengchengyuan]{有限生成投射生成元}[progenerator]。投射性指沿任意满同态的提升性质，生成元性质采用\cref{b2:thm:generator-characterizations}的任一等价条件。
\end{definition}

\cref{b2:prop:finite-projective-summand}与生成元判据合在一起说明，这等价于同时存在有限整数 \(m,n\) 及模 \(Q,Q'\)，使
\[
P\oplus Q\cong R^m,\qquad R\oplus Q'\cong P^n.
\]
第一个分解把 \(P\) 放在有限自由模中作为直和因子，控制其有限生成与投射性；第二个分解把正则模放在有限份 \(P\) 中，保证它能够产生全部模。这两类分解不能互相替代，\cref{b2:prop:progenerator-hypothesis-counterexamples}将分别检验三个条件缺失时的情形。

\ProseParagraphBreak
有限生成性用一组元素来表述，但范畴等价并不指定两个对象底集合的元素如何对应。要在等价下追踪投射模的有限生成性，需要把它改写成能由 Hom 和直和识别的条件：从这个模到任意直和的同态，是否总能只用有限个目标分量表达？

\begin{definition}\label{b2:def:compact-module-object}
设 \(R\) 为含幺环，\(P\) 为幺左 \(R\) 模。对任意指标集 \(I\) 及幺左 \(R\) 模族 \((M_i)_{i\in I}\)，若由各分量包含映射诱导的自然映射
\[
\bigoplus_{i\in I}\operatorname{Hom}_R(P,M_i)
\longrightarrow\operatorname{Hom}_R\left(P,\bigoplus_{i\in I}M_i\right)
\]
均为双射，则称 \(P\) 为左 \(R\) 模范畴中的\term[jin-duixiang]{紧对象}[compact object]。此处专指 Hom 保持任意直和的范畴条件，不指拓扑上的紧性，也不同于以 Hom 保持滤过余极限刻画的有限展示性。
\end{definition}

\begin{proposition}\label{b2:prop:compact-projective-finite}
设 \(R\) 为含幺环，所论模均为幺左 \(R\) 模。每个有限生成左 \(R\) 模都是左 \(R\) 模范畴中的紧对象。若左 \(R\) 模 \(P\) 还投射，则 \(P\) 有限生成，当且仅当 \(P\) 是紧对象。
\end{proposition}
\begin{proof}
自然映射始终单射：若有限支集族 \((f_i)_i\) 合成到直和中的同态为零，再接上第 \(j\) 个坐标投影，便得到 \(f_j=0\)，对每个 \(j\) 都成立。若 \(P\) 有有限生成组，任意 \(f:P\rightarrow\bigoplus_iM_i\) 在这些生成元上的像共同只有有限支集，故全像落在一个有限子直和中。于是 \(f\) 是有限多个分量同态之和，证明满射。

反过来，设 \(P\) 投射且紧。由\cref{b2:thm:projective-characterizations}，可取分裂自由满射 \(q:R^{(I)}\rightarrow P\) 及截面 \(s:P\rightarrow R^{(I)}\)。紧性使 \(s\) 的像落在某个有限子直和 \(R^J\) 中。限制 \(q\) 到 \(R^J\) 后仍满足 \(qs=\operatorname{id}_P\)，所以 \(P\) 为有限自由模的直和因子，因而有限生成。
\end{proof}

对投射对象，这个判据把有限生成性转成了由 Hom 与直和描述的条件，后面的等价定理将使用这一形式。

\subsection{矩阵环与原环的模范畴}
\label{b2:subsec:1601:03}

\begin{proposition}\label{b2:prop:matrix-morita-equivalence}
设 \(R\) 为含幺环、\(n\ge1\) 为整数，所论模均为幺性左模。令 \(A=M_n(R)\)、\(e=E_{11}\)，将 \(eAe\) 按 \(r\mapsto rE_{11}\) 识别为 \(R\)。以下函子互为拟逆：
\[
\mathcal L:R\text{-}\mathbf{Mod}\longrightarrow A\text{-}\mathbf{Mod},
\quad N\longmapsto N^n,
\qquad
\mathcal K:A\text{-}\mathbf{Mod}\longrightarrow R\text{-}\mathbf{Mod},
\quad M\longmapsto eM.
\]
\(\mathcal L\) 使用矩阵乘列的左作用，\(\mathcal K\) 使用角环作用；在同态上分别逐坐标作用和限制到 \(eM\)。
\end{proposition}
\begin{proof}
\(eN^n\) 恰为只有第一坐标可能非零的列，取该坐标自然同构到 \(N\)。另一方向用 \(E_{1j}\) 将第 \(j\) 个位置的数据读到第一角，再用 \(E_{i1}\) 写回第 \(i\) 个位置。相应地，定义
\[
\Theta_M:(eM)^n\longrightarrow M,\qquad
(m_1,\ldots,m_n)\longmapsto\sum_iE_{i1}m_i,
\]
并令 \(\Psi_M(m)=(E_{11}m,E_{12}m,\ldots,E_{1n}m)\)。每个坐标都在 \(eM\)。由 \(E_{1j}E_{i1}=\delta_{ij}e\) 与 \(\sum_iE_{i1}E_{1i}=1\)，两映射互逆。

对 \(a=(a_{ij})\in A\)，恒等式 \(aE_{j1}=\sum_iE_{i1}(a_{ij}E_{11})\) 说明 \(\Theta_M\) 把列上的矩阵作用变为 \(M\) 的原作用，故 \(A\) 线性。两映射只使用矩阵单位的作用，与任意 \(A\) 模同态交换，所以构成自然同构。第一坐标同构也显然与 \(R\) 模同态相容，因此两函子确为范畴等价。
\end{proof}

这个等价不要求 \(R\) 交换。例如 \(M_n(R)\) 的正则左模被 \(\mathcal K\) 送到首行模 \(eA\cong R^n\)。取 \(n=2\)，矩阵 \(m=\begin{psmallmatrix}a&b\\c&d\end{psmallmatrix}\) 对应于
\[
\Psi_A(m)=\left(
\begin{pmatrix}a&b\\0&0\end{pmatrix},
\begin{pmatrix}c&d\\0&0\end{pmatrix}
\right).
\]
逆映射把第一个首行留在第一行，再由 \(E_{21}\) 把第二个首行移到第二行，便恢复 \(m\)。矩阵的全部条目因此可以从首行模的两份副本恢复，而 \(eA\) 本身是秩二的 \(R\) 模。这里正则模 \({}_AA\) 的对应对象是 \({}_RR^n\)，等价不要求两边指定的正则模互相对应。取出 \(eM\)，再由矩阵单位恢复整个 \(M\)，给出了后面用双模张量积重构模的具体模型。

\begin{proposition}\label{b2:prop:trace-generator-criterion}
设 \(R\) 为含幺环、\(P\) 为幺左 \(R\) 模。子集
\[
\operatorname{tr}_R(P)=\sum_{f\in\operatorname{Hom}_R(P,R)}f(P)
\]
是 \(R\) 的双边理想，称为 \(P\) 的\term[ji-lixiang]{迹理想}[trace ideal]。以下条件等价：\(P\) 是全部幺左 \(R\) 模范畴的生成元；\(\operatorname{tr}_R(P)=R\)；存在非负整数 \(n\)、元素 \(p_1,\ldots,p_n\in P\) 和左 \(R\) 线性映射 \(f_1,\ldots,f_n:P\to R\)，使
\[
\sum_{i=1}^n f_i(p_i)=1_R.
\]
这里 \(n=0\) 时按空和为零理解，不要求 \(P\) 有限生成或投射。
\end{proposition}
\symidx{\(\operatorname{tr}_R(P)\)：模的迹理想}
\begin{proof}
每个 \(f(P)\) 都是正则左模的子模，所以它们的和是左理想。对任意 \(r\in R\)，映射 \(f_r:P\to R\)、\(f_r(p)=f(p)r\) 仍左线性，因为 \(f_r(ap)=af(p)r=af_r(p)\)。故迹理想也对右乘封闭，是双边理想。

迹理想为全环，当且仅当其中含有一；而这个和中的每个元素都由有限多个像相加得到，所以这又等价于所述有限和等式。若等式成立，定义
\[
q:P^n\longrightarrow R,\qquad (x_i)_i\longmapsto\sum_i f_i(x_i),
\qquad
s:R\longrightarrow P^n,\qquad r\longmapsto(rp_i)_i.
\]
两映射都左线性，且 \(q(s(r))=r\sum_i f_i(p_i)=r\)，故 \(R\) 是 \(P^n\) 的直和因子，由\cref{b2:thm:generator-characterizations}，\(P\) 为生成元。反过来，生成元判据给出有限份 \(P\) 到 \(R\) 的分裂满射。取一的原像 \((p_i)_i\)，并将满射限制到第 \(i\) 个分量得到 \(f_i\)，就有 \(\sum_i f_i(p_i)=1_R\)。
\end{proof}

迹判据中的有限和控制的是 \(1_R\)。相比之下，\cref{b2:prop:finite-dual-basis}中的有限对偶基控制的是每个 \(p\in P\) 的重构式，因而控制 \(\operatorname{id}_P\)。这两种有限数据分别检验生成元性质与有限生成投射性。

\begin{proposition}\label{b2:prop:object-detection-not-generator}
设 \(k\) 为域、\(R=k[\varepsilon]/(\varepsilon^2)\)，\(\varepsilon\) 的剩余类仍记为 \(\varepsilon\)，并令 \(P=R/(\varepsilon)\) 为幺左 \(R\) 模。对每个非零幺左 \(R\) 模 \(M\)，都有 \(\operatorname{Hom}_R(P,M)\ne0\)；但 \(\operatorname{Hom}_R(P,-)\) 不忠实，\(P\) 不是生成元。
\end{proposition}
\begin{proof}
在 \(M\) 中取 \(m\ne0\)。若 \(\varepsilon m=0\)，令 \(v=m\)；否则令 \(v=\varepsilon m\)，由 \(\varepsilon^2=0\) 仍有 \(v\ne0\) 且 \(\varepsilon v=0\)。于是 \(r+(\varepsilon)\mapsto rv\) 良定义，给出 \(P\to M\) 的非零左模同态。

考虑非零左线性映射 \(u:R\to R\)、\(u(r)=r\varepsilon\)。每个 \(h:P\to R\) 的像都被 \(\varepsilon\) 零化；写 \(r=a+b\varepsilon\) 可知 \(\varepsilon r=a\varepsilon=0\) 当且仅当 \(r\in k\varepsilon\)，故 \(h(P)\subseteq k\varepsilon\)，从而 \(uh=0\)。于是 \(u\ne0\) 和零同态经 Hom 函子后成为同一个映射，忠实性失败，\(P\) 不是生成元。
\end{proof}

\begin{proposition}\label{b2:prop:progenerator-hypothesis-counterexamples}
设 \(k\) 为域，所论模均为幺性左模。有限生成投射生成元定义中的三个条件，均不能由另外两个推出：
\begin{enumerate}
\item 在 \(R=k\times k\) 上，\(P=k\times0\) 有限生成且投射，但不是生成元。
\item 在 \(R=\mathbb Z\) 上，\(P=\mathbb Z\oplus\mathbb Z/2\mathbb Z\) 有限生成且为生成元，但不投射；\(\operatorname{Hom}_{\mathbb Z}(P,-)\) 不保持所有满同态。
\item 在 \(R=k\) 上，\(P=\bigoplus_{i\ge1}ku_i\) 投射且为生成元，但不有限生成，也不是紧对象；\(\operatorname{Hom}_k(P,-)\) 不保持任意直和。
\end{enumerate}
\end{proposition}
\begin{proof}
在第一例中，令 \(e=(1,0)\)。\(P=Re\) 是 \(R=Re\oplus R(1-e)\) 的直和因子，又由 \(e\) 生成，所以有限生成且投射。非零模 \(Q=0\times k\) 满足 \(eQ=0\)。若 \(h:P\to Q\)，则 \(h(e)=h(e^2)=eh(e)=0\)，所以 \(h=0\)。故 \(Q\) 上的恒等映射与零同态无法由 \(P\) 检测，\(P\) 不是生成元。

第二例的两个直和分量各由一个元素生成，而 \(\mathbb Z\) 是 \(P\) 的直和因子，所以由生成元判据，\(P\) 是有限生成的生成元。取自然满射 \(q:\mathbb Z\to\mathbb Z/2\mathbb Z\) 及投影 \(f:P\to\mathbb Z/2\mathbb Z\)。若有 \(h:P\to\mathbb Z\) 满足 \(qh=f\)，则 \(h\) 在 \(\mathbb Z/2\mathbb Z\) 分量上的限制应是从二阶元素到整数的同态，只能为零；但 \(f\) 在此分量上为恒等，矛盾。因此 \(P\) 不投射，\(\operatorname{Hom}_{\mathbb Z}(P,q)\) 也不满射。

第三例中，\(P\) 自由，因而投射；它含 \(ku_1\cong k\) 为直和因子，故是生成元。有限个向量的支集之并仍有限，所以不能生成全部 \(u_i\)，\(P\) 不有限生成。将恒等映射看成 \(P\to\bigoplus_{i\ge1}ku_i\)。如果它来自 \(\bigoplus_i\operatorname{Hom}_k(P,ku_i)\)，其像就会落在有限个坐标的直和中，不能包含所有 \(u_i\)，矛盾。所以 \(P\) 不是紧对象，Hom 函子不能保持这个直和。
\end{proof}

% !TeX root = ../../Book2.tex
% 纲要标题：### 16.2 双模实现的函子
% 纲要位置：计划纲要.md，第 1570 行。

\section{双模实现的函子}
\label{b2:sec:1602}

为了把两个模范畴联系起来，需要一个同时接受两侧标量作用的模。先从任意 \((R,S)\) 双模 \(P\) 出发，写出张量函子、Hom 函子及它们的伴随映射；然后再把 \(S\) 选为 \(P\) 的自同态环的反环。这样可以分别检查构造本身的合法性，以及何种条件使它成为等价。

% 纲要要点（供目录核对）：
%
% * 张量函子与 Hom 函子；调用第7.1节的有限投射 Hom–张量同构
% * 自同态环和反环约定；逐式核对预复合和右作用，正式证明任意环上的行列互逆双模及一般模上的余单位
% * 伴随的单位及余单位
%

\subsection{张量函子与 Hom 函子}
\label{b2:subsec:1602:01}

设 \(R,S\) 为含幺环，给定幺性双模 \({}_RP_S\)，定义
\[
F(M)=\operatorname{Hom}_R(P,M),\qquad
G(N)=P\otimes_SN.
\]
这里 \(M\) 为幺左 \(R\) 模，\(N\) 为幺左 \(S\) 模，因此 \(F:R\text{-}\mathbf{Mod}\to S\text{-}\mathbf{Mod}\)、\(G:S\text{-}\mathbf{Mod}\to R\text{-}\mathbf{Mod}\)。\(G(N)\) 的左 \(R\) 作用留在第一个因子；\(F(M)\) 的左 \(S\) 作用由
\begin{equation}\label{b2:eq:morita-hom-action}
(sf)(p)=f(ps)
\end{equation}
给出：\(s\) 先在输入端作用，再由 \(f\) 求值，因而这是预复合。它仍为左 \(R\) 线性映射，因为 \(P\) 的两侧作用相容；而
\[
\bigl(s_1(s_2f)\bigr)(p)=f\bigl((ps_1)s_2\bigr)
=\bigl((s_1s_2)f\bigr)(p),
\]
所以确为左 \(S\) 作用。

\ProseParagraphBreak
对 \(a:M\rightarrow M'\)，令 \(F(a)(f)=af\)；对 \(b:N\rightarrow N'\)，令 \(G(b)=\operatorname{id}_P\otimes b\)。复合与恒等关系由同态复合和张量函子性得到。它们都是加性函子。由\cref{b2:thm:hom-tensor-adjunction}，有关于两变量自然的对应
\[
\operatorname{Hom}_R(P\otimes_SN,M)
\cong\operatorname{Hom}_S\bigl(N,\operatorname{Hom}_R(P,M)\bigr),
\]
把 \(h\) 送到 \(n\mapsto\bigl(p\mapsto h(p\otimes n)\bigr)\)。因此 \(G\) 左伴随于 \(F\)。一般地，\(F\) 左正合，\(G\) 右正合；\(P\) 左投射时，\cref{b2:cor:projective-hom-exact}才使 \(F\) 正合。

\ProseParagraphBreak
若这里的双模 \({}_RP_S\) 作为左 \(R\) 模有限生成并投射，\cref{b2:prop:finite-projective-hom-tensor}便把已经定义的 \(F(M)=\operatorname{Hom}_R(P,M)\) 写成张量函子。具体地，令 \(P^*=\operatorname{Hom}_R(P,R)\)，它接受 \((S,R)\) 双模作用
\[
(s\lambda)(p)=\lambda(ps),\qquad
(\lambda r)(p)=\lambda(p)r.
\]
该命题给出关于每个左 \(R\) 模 \(M\) 自然的左 \(S\) 模同构
\[
P^*\otimes_RM\longrightarrow F(M),\qquad
\lambda\otimes m\longmapsto\bigl(p\mapsto\lambda(p)m\bigr).
\]
目标的左 \(S\) 作用正是式\eqref{b2:eq:morita-hom-action}中的作用，因此这个同构保留当前的双模侧别。公式中的标量 \(\lambda(p)\) 作用在 \(m\) 左侧；交换这两个符号在一般环上并没有意义。

\subsection{自同态环和反环约定}
\label{b2:subsec:1602:02}

现在固定一个含幺环 \(R\) 及幺左 \(R\) 模 \(P\)，取
\[
\Delta=\operatorname{End}_R(P),\qquad S=\Delta^{\mathrm{op}}.
\]
\symidx{\(S=\operatorname{End}_R(P)^{\mathrm{op}}\)：投射生成元对应的自同态反环}
若 \(s\in S\) 对应自同态 \(\alpha_s\in\Delta\)，定义 \(ps=\alpha_s(p)\)。因为
\[
\alpha_{st}=\alpha_t\circ\alpha_s,
\qquad (ps)t=\alpha_t\bigl(\alpha_s(p)\bigr)=p(st),
\]
这是右 \(S\) 作用；各 \(\alpha_s\) 左 \(R\) 线性，保证它与左作用相容。因此 \(P\) 成为 \((R,S)\) 双模，上一个小节的函子与伴随立即适用。

两侧作用可记为
\[
{}_RP_S,\qquad {}_SP^*_R,\qquad
P^*=\operatorname{Hom}_R(P,R),\qquad
S=\operatorname{End}_R(P)^{\mathrm{op}}.
\]
自同态的复合在右 \(S\) 作用中反序一次，因而这里必须取反环。

\ProseParagraphBreak
在此约定下，\(F(P)=\operatorname{End}_R(P)\) 作为左 \(S\) 模自然同构于正则左模 \(S\)：将 \(\alpha_s\) 对应于 \(s\)。事实上，\(t\alpha_s\) 的作用是先做 \(\alpha_t\) 再做 \(\alpha_s\)，即 \(\alpha_s\alpha_t=\alpha_{ts}\)，正对应左乘 \(t\)。所以在将来由这些函子给出的等价中，与正则左 \(S\) 模对应的是选定的 \(P\)，未必是正则左 \(R\) 模。这是一个左模识别；作为环，\(\Delta=S^{\mathrm{op}}\)，不能因此省去反环。

\ProseParagraphBreak
最基本的例子是 \(P={}_RR\)。每个左线性自同态都是右乘某个 \(a\in R\)，所以 \(\Delta\cong R^{\mathrm{op}}\)、\(S\cong R\)，Hom 在一处求值后恢复原模。若取左自由模 \(P=R^{1\times n}\)、\(n\ge1\)，写成行向量方便记录右作用；非交换系数不能随意搬到另一侧，所以不能把以下公式直接改成左乘列：每个左线性自同态唯一为右乘矩阵 \(B\in M_n(R)\)，两次右乘的算子复合满足
\[
D_BD_C(p)=(pC)B=p(CB).
\]
所以 \(S\cong M_n(R)\)。映射 \(f:P\rightarrow M\) 对应其在各标准行上的值组成的列
\[
\bigl(f(e_1),\ldots,f(e_n)\bigr)^{\mathsf T}\in M^n.
\]
由 \(e_iB=\sum_j b_{ij}e_j\)，左 \(S\) 作用恰为矩阵乘这个列。这便与\cref{b2:prop:matrix-morita-equivalence}的显式等价相合。

\subsection{伴随的单位及余单位}
\label{b2:subsec:1602:03}

对任意 \((R,S)\) 双模，伴随的单位和余单位由\cref{b2:thm:hom-tensor-adjunction}的两种求值公式给出：
\begin{align}
\eta_N:N&\longrightarrow FG(N),
&\eta_N(n)(p)&=p\otimes n,\label{b2:eq:morita-unit}\\
\varepsilon_M:GF(M)&\longrightarrow M,
&\varepsilon_M(p\otimes f)&=f(p).\label{b2:eq:morita-counit}
\end{align}
第一式比较先取张量再取 Hom 所得对象与原来的 \(N\)，第二式则将一个输入 \(p\) 与同态 \(f\) 配对，尝试恢复原来的 \(M\)。第二式的平衡性为 \(f(ps)=(sf)(p)\)，第一式的左 \(S\) 线性为 \(p\otimes sn=ps\otimes n\)。相应的左 \(R\) 线性分别来自 \(P\) 的左作用和 \(f\) 的线性。因此两者都是所在模范畴中的同态，不只是集合映射。

\begin{proposition}\label{b2:prop:morita-adjunction-maps}
设 \(R,S\) 为含幺环，\({}_RP_S\) 为幺性 \((R,S)\) 双模，令 \(F=\operatorname{Hom}_R(P,-)\)、\(G=P\otimes_S-\)。对左 \(S\) 模 \(N\) 和左 \(R\) 模 \(M\)，定义
\[
\eta_N:N\longrightarrow FG(N),\quad n\longmapsto\bigl(p\longmapsto p\otimes n\bigr),
\qquad
\varepsilon_M:GF(M)\longrightarrow M,\quad p\otimes f\longmapsto f(p).
\]
这些映射关于 \(N,M\) 自然，并满足
\begin{equation}\label{b2:eq:morita-triangles}
\varepsilon_{G(N)}G(\eta_N)=\operatorname{id}_{G(N)},\qquad
F(\varepsilon_M)\eta_{F(M)}=\operatorname{id}_{F(M)}.
\end{equation}
若 \(S=\operatorname{End}_R(P)^{\mathrm{op}}\)，且 \(P\) 的右 \(S\) 作用取自其左 \(R\) 线性自同态，则 \(\eta_S\) 与 \(\varepsilon_P\) 都是同构，不要求 \(P\) 已是生成元。
\end{proposition}
\begin{proof}
若 \(b:N\rightarrow N'\)，两种从 \(N\) 到 \(FG(N')\) 的路径在 \(n,p\) 处都为 \(p\otimes b(n)\)；若 \(a:M\rightarrow M'\)，余单位的两种路径在 \(p\otimes f\) 处都为 \(a\bigl(f(p)\bigr)\)，所以自然。第一条三角恒等式把 \(p\otimes n\) 送到 \(\eta_N(n)(p)=p\otimes n\)；第二条把 \(f\) 送到函数 \(p\mapsto\varepsilon_M(p\otimes f)=f(p)\)，故都成立。

在所给自同态环约定下，识别 \(P\otimes_SS\cong P\) 及 \(F(P)\cong S\)。\(\eta_S(s)\) 对应 \(p\mapsto ps=\alpha_s(p)\)，再对应到 \(s\)，所以是恒等映射。\(\varepsilon_P\) 在这些识别下为 \(P\otimes_SS\rightarrow P\)、\(p\otimes s\mapsto ps\)，是张量积的单位同构。
\end{proof}

这些特殊分量为同构，并不说明其他模上的分量也是同构。为了在全部模上得到同构，需要把模写成适当直和之间映射的余核，并证明函子保留相应的直和与正合尾段。有限生成、投射和生成元三个条件将在下一节分别承担这些任务。

\begin{proposition}[行列双模的互逆配对]\label{b2:prop:matrix-row-column-bimodules}
设 \(R\) 为含幺环、\(n\ge1\)，令 \(S=M_n(R)\)、\({}_RP_S=R^{1\times n}\)、\({}_SQ_R=R^{n\times1}\)，各作用由矩阵乘法给出。令 \(u_i\)、\(v_i\) 分别为第 \(i\) 个标准行、标准列，则矩阵乘法诱导双模同构
\[
P\otimes_SQ\longrightarrow R,\quad p\otimes q\longmapsto pq,
\qquad Q\otimes_RP\longrightarrow S,\quad q\otimes p\longmapsto qp.
\]
它们的逆分别为 \(r\mapsto ru_1\otimes v_1\) 和
\[
(b_{ij})\longmapsto\sum_{i,j}v_i b_{ij}\otimes u_j.
\]
对任意幺左 \(R\) 模 \(M\)，将 \(\operatorname{Hom}_R(P,M)\) 识别为左 \(S\) 模 \(M^n\) 后，余单位也是同构：
\[
P\otimes_SM^n\longrightarrow M,\quad
(p_i)_i\otimes(m_i)_i^{\mathsf T}\longmapsto\sum_i p_i m_i,
\]
其逆为 \(m\mapsto u_1\otimes(m,0,\ldots,0)^{\mathsf T}\)。
\end{proposition}
\begin{proof}
结合律给出两种矩阵配对的平衡性与外侧作用相容性。为核对第一同构的逆，将行列按标准基展开，得到
\[
(p_i u_i)\otimes(v_jq_j)
=u_1\otimes (p_iE_{1i})(v_jq_j)
=\delta_{ij}(p_iq_j)u_1\otimes v_1.
\]
第一步在 \(S\) 上移动 \(p_iE_{1i}\)，最后一步移动 \((p_iq_j)E_{11}\)。因此全张量等于 \((pq)u_1\otimes v_1\)，两复合均为恒等。第二同构的逆把 \(E_{ij}\) 送到 \(v_i\otimes u_j\)；对任意纯张量，\(R\) 上的平衡关系给出
\[
q\otimes p=\sum_{i,j}v_i(q_i p_j)\otimes u_j,
\]
恰为外积矩阵 \(qp\) 的逆像。

最后，任意行 \(p\) 可写成 \(u_1B\)，其中 \(B\) 的首行为 \(p\)、其余行为零。在张量中把 \(B\) 移到 \(M^n\) 上，得到
\[
p\otimes(m_i)_i^{\mathsf T}
=u_1\otimes\left(\sum_i p_i m_i,0,\ldots,0\right)^{\mathsf T}.
\]
这就验证了所列余单位及其逆在任意 \(M\) 上的两个复合。所有系数保持原次序，证明不要求 \(R\) 交换。
\end{proof}

% !TeX root = ../../Book2.tex
% 纲要标题：### 16.3 Morita 等价定理
% 纲要位置：计划纲要.md，第 1576 行。

\section{Morita 等价定理}
\label{b2:sec:1603}

第五章把 Morita 等价定义为模范畴的等价，却没有说明怎样从环本身判断它。本节证明：选择一个有限生成投射生成元，再取其自同态环的反环，恰好得到全部这样的等价。先从一个已给的范畴等价恢复这个模，再证明投射生成元确实使上一节的单位、余单位在每个对象上成为同构。

% 纲要要点（供目录核对）：
%
% * 从范畴等价恢复投射生成元；由双积恢复加法，先保持紧性和投射性，再恢复有限生成性
% * 由投射生成元构造等价；将单位与余单位从生成对象经任意直和、展示及余核推广到全部模
% * 互逆双模及右模版本
% * 满幂等元与角环的 Morita 等价；正式纳入非零对象检测判据和求值逆，并证明双边理想及商环对应
%

\subsection{从范畴等价恢复投射生成元}
\label{b2:subsec:1603:01}

\begin{lemma}\label{b2:lem:module-equivalence-structure}
设 \(R,S\) 为含幺环，\(H:R\text{-}\mathbf{Mod}\rightarrow S\text{-}\mathbf{Mod}\) 为幺左模范畴之间的等价，则 \(H\) 为加性函子，保持并反映核、余核、正合列、投射性与内射性，也保持任意直和，并将生成元送到生成元。
\end{lemma}
\begin{proof}
核、余核及直和由泛性质刻画，投射性、内射性与生成元也可用同态的提升、延拓和商对象刻画，因此这些性质都能随等价转移。需要先说明的是，同态的加法本身也能由范畴结构恢复。由\cref{b2:prop:equivalence-preserves-universal-objects}，等价保持零对象及任意积、余积，特别保持任意直和。零同态是经过零对象的复合，因此也被保持。对同态 \(f,g:M\rightarrow N\)，加法可写成
\[
f+g=\nabla_N(f\oplus g)\Delta_M,
\]
其中 \(\Delta_M:M\rightarrow M\oplus M\) 的两个投影均为恒等，\(\nabla_N:N\oplus N\rightarrow N\) 的两个限制均为恒等。直和同时是积与余积，故等价在相应的自然直和识别下保持这两个映射及此式，所以保持加法。无需在模范畴的等价之外另加一个加性要求。

等价保持并反映单态射和满态射的消去性质；模范畴中的单态射就是单射，因为非零核元素可以通过从 \(R\) 出发的同态检测。满态射就是满射，因为若像不为全体，非零商映射与零映射在其像上相同，会破坏右消去性。因此单射与满射也被保持并反映。

以核为例，设 \(j:K\rightarrow M\) 是 \(f:M\rightarrow N\) 的核，等价记为 \(H\)。若 \(a:Y\rightarrow H(M)\) 满足 \(H(f)a=0\)，取同构 \(u:H(X)\rightarrow Y\)，再由全忠实性把 \(au\) 唯一提升为 \(b:X\rightarrow M\)。忠实性给出 \(fb=0\)，故 \(b=jc\) 且 \(c\) 唯一。于是 \(a=H(j)H(c)u^{-1}\)，而任何别的分解提升后都等于 \(c\)。所以 \(H(j)\) 是核。

对余核 \(q:N\rightarrow C\)，给定 \(a:H(N)\rightarrow Y\) 且 \(aH(f)=0\)，选 \(u:H(X)\rightarrow Y\)，把 \(u^{-1}a\) 提升为 \(b:N\rightarrow X\)，则 \(bf=0\)，由余核性质唯一写成 \(b=cq\)，所以 \(a=uH(c)H(q)\) 且分解唯一。故余核也保持。像可表为余核的核，因此也保持。对拟逆使用同样论证，得到反映性，正合列因核与像的相等而双向保持。

投射性由沿满射提升的性质刻画。将一个提升问题经拟逆送回，满射仍满，再作提升并经等价送出，就保持投射性；反向相同。对偶地，内射性由沿单射延拓的性质刻画；等价及拟逆都保持单射，故也保持并反映内射性。若 \(P\) 是生成元，每个源范畴对象都为 \(P\) 某个直和的商。等价保持直和与满射，且每个目标对象都同构于某个 \(H(X)\)，故每个目标对象都为 \(H(P)\) 某个直和的商。因此 \(H(P)\) 也是生成元。
\end{proof}

目标范畴的正则模 \({}_SS\) 是投射生成元，而且 \(\operatorname{Hom}_S(S,Y)\) 在 \(1\) 处求值就恢复 \(Y\)。因此把它经拟逆送回得到 \(P=\Psi(S)\)，应当能用 \(\operatorname{Hom}_R(P,-)\) 恢复原等价。投射性与生成元性质已有范畴刻画；有限生成性则不能直接用元素和生成元的个数转移，需要先把正则模的紧性转移给 \(P\)，再利用 \(P\) 的投射性。

\begin{theorem}[由范畴等价恢复投射生成元]\label{b2:thm:morita-converse}
设 \(R,S\) 为含幺环，\(\Phi:R\text{-}\mathbf{Mod}\rightarrow S\text{-}\mathbf{Mod}\) 为幺左模范畴的等价，取其拟逆 \(\Psi\)，令 \(P=\Psi({}_SS)\)，则幺左 \(R\) 模 \(P\) 为有限生成投射生成元，并有环同构
\[
S\cong\operatorname{End}_R(P)^{\mathrm{op}}.
\]
通过这个同构赋予 Hom 左 \(S\) 作用后，\(\Phi\) 自然同构于 \(\operatorname{Hom}_R(P,-)\)。
\end{theorem}
\begin{proof}
正则模 \(S\) 投射且为生成元，\cref{b2:lem:module-equivalence-structure}使 \(P=\Psi(S)\) 也投射且为生成元。固定拟逆提供的自然同构 \(\nu:\operatorname{Id}_{S\text{-}\mathbf{Mod}}\rightarrow\Phi\Psi\)。为证明有限生成，先证明 \(P\) 紧。对任意幺左 \(R\) 模族 \((M_i)\)，有自然同构
\[
\begin{aligned}
\operatorname{Hom}_R\left(P,\bigoplus_iM_i\right)
&\cong\operatorname{Hom}_S\left(S,\Phi\left(\bigoplus_iM_i\right)\right)\\
&\cong\operatorname{Hom}_S\left(S,\bigoplus_i\Phi(M_i)\right)\\
&\cong\bigoplus_i\operatorname{Hom}_S\bigl(S,\Phi(M_i)\bigr)\\
&\cong\bigoplus_i\operatorname{Hom}_R(P,M_i).
\end{aligned}
\]
第一项同构来自 \(\Phi\) 的全忠实性及 \(\nu_S:S\cong\Phi(P)\)，第二项来自等价保持直和。第三项使用正则模 \(S\) 由 \(1\) 生成：从 \(S\) 到直和的同态由 \(1\) 的像决定，而这个像只有有限支集。最后一项再用全忠实性及 \(\nu_S\)。这些同构与各直和嵌入相容，所以其复合正是紧性中的自然映射之逆。\cref{b2:prop:compact-projective-finite}于是说明投射模 \(P\) 有限生成。

等价加性且全忠实，故 \(\Psi\) 在自同态环上给出环同构
\[
\operatorname{End}_S({}_SS)\cong\operatorname{End}_R(P).
\]
由\cref{b2:prop:regular-hom-endomorphism}，正则左模的自同态是右乘映射 \(r_s:x\mapsto xs\)，并有 \(r_s\circ r_t=r_{ts}\)。因此 \(s\mapsto\Psi(r_s)\) 给出 \(S\cong\operatorname{End}_R(P)^{\mathrm{op}}\)，相应的右作用为 \(ps=\Psi(r_s)(p)\)。

对每个 \(M\)，再用正则模的求值对应
\[
\operatorname{Hom}_R(P,M)
\cong\operatorname{Hom}_S(S,\Phi(M))
\cong\Phi(M).
\]
具体地，\(h:P\rightarrow M\) 先对应 \(a=\Phi(h)\nu_S:S\rightarrow\Phi(M)\)，再对应 \(a(1)\)。后复合 \(M\rightarrow M'\) 的同态与这两步相容，故所得同构对 \(M\) 自然。它也保持左 \(S\) 作用：Hom 中的 \(s\cdot h\) 为 \(h\circ\Psi(r_s)\)，而 \(\nu\) 的自然性给出
\[
\Phi\bigl(h\circ\Psi(r_s)\bigr)\nu_S
=\Phi(h)\nu_S r_s=a r_s.
\]
在 \(1\) 处求值，右边为 \(a(s)=s\cdot a(1)\)，所以该自然同构为左 \(S\) 模同构。
\end{proof}

\subsection{投射生成元与等价构造}
\label{b2:subsec:1603:02}

\begin{theorem}[Morita 等价]\label{b2:thm:morita-equivalence}
设 \(R\) 为含幺环，\(P\) 为幺左 \(R\) 模范畴的有限生成投射生成元，令 \(S=\operatorname{End}_R(P)^{\mathrm{op}}\)，并使 \(s\in S\) 在 \(P\) 上的右作用为 \(ps=\alpha_s(p)\)，其中 \(\alpha_s\) 是 \(s\) 对应的左 \(R\) 线性自同态，则
\[
F=\operatorname{Hom}_R(P,-),\qquad G=P\otimes_S-
\]
分别为从幺左 \(R\) 模到幺左 \(S\) 模、从幺左 \(S\) 模到幺左 \(R\) 模的互逆等价，而且式\eqref{b2:eq:morita-unit}与式\eqref{b2:eq:morita-counit}的单位、余单位在每个模上都是同构。结合\cref{b2:thm:morita-converse}，任意两个含幺环 \(R,S\) 在 Morita 意义下等价，当且仅当存在幺左 \(R\) 模的有限生成投射生成元 \(P\) 使 \(S\cong\operatorname{End}_R(P)^{\mathrm{op}}\)。
\end{theorem}
\begin{proof}
上一节已知 \(\varepsilon_P\) 与 \(\eta_S\) 为同构。要把它们扩展到所有对象，可以先扩展到 \(P\) 或 \(S\) 的任意直和，再把任意模写成这些直和之间某个映射的余核。这里有限生成性与\cref{b2:prop:compact-projective-finite}保证 \(F\) 保持任意直和；投射性与\cref{b2:cor:projective-hom-exact}保证 \(F\) 正合。\cref{b2:prop:tensor-direct-sums}与\cref{b2:thm:tensor-right-exact}说明 \(G\) 保持直和且右正合。因此 \(FG,GF\) 都保持直和与正合尾段。

由\cref{b2:prop:morita-adjunction-maps}，\(\varepsilon_P\) 为同构；因两边函子都保持直和，且求值公式与各嵌入相容，每个 \(\varepsilon_{P^{(I)}}\) 也是同构。对任意幺左 \(R\) 模 \(M\)，生成元性质给出满射 \(q:P^{(J)}\rightarrow M\)。记其核为 \(K\)，再对 \(K\) 使用生成元性质，得到满射 \(P^{(I)}\rightarrow K\)。将它与核的嵌入复合，得到 \(d:P^{(I)}\rightarrow P^{(J)}\) 及正合列
\[
P^{(I)}\longrightarrow P^{(J)}\longrightarrow M\longrightarrow0.
\]
生成元性质的这两次使用，分别保证 \(M\) 有由 \(P\) 给出的生成元和关系。应用右正合函子 \(GF\)，再用余单位的自然性，得到交换的正合尾段
\[
\begin{tikzcd}[column sep=2.6em,row sep=2.5em]
GF(P^{(I)}) \arrow[r,"GF(d)"] \arrow[d,"\varepsilon_{P^{(I)}}"'] &
GF(P^{(J)}) \arrow[r,"GF(q)"] \arrow[d,"\varepsilon_{P^{(J)}}"'] &
GF(M) \arrow[r] \arrow[d,"\varepsilon_M"'] & 0\\
P^{(I)} \arrow[r,"d"'] &
P^{(J)} \arrow[r,"q"'] &
M \arrow[r] & 0
\end{tikzcd}
\]
前两列的竖直映射为同构，并把 \(\operatorname{im}GF(d)\) 对应到 \(\operatorname{im}d\)，因而在对这两个像取商后诱导同构。两行的最后一个对象都是相应第一映射的余核，故由余核泛性质，这个同构正是 \(\varepsilon_M\)。

已知 \(\eta_S\) 为同构，两边保持直和，使每个 \(\eta_{S^{(I)}}\) 也为同构。任意幺左 \(S\) 模 \(N\) 本来就有自由展示
\[
S^{(I)}\longrightarrow S^{(J)}\longrightarrow N\longrightarrow0,
\]
这里指标集允许无限；这一侧用的是正则模 \(S\) 的生成元性质，不再需要 \(P\) 生成 \(R\) 模。把前两个映射记为 \(d':S^{(I)}\to S^{(J)}\)、\(q':S^{(J)}\to N\)。应用右正合函子 \(FG\)，由单位的自然性得到
\[
\begin{tikzcd}[column sep=2.6em,row sep=2.5em]
S^{(I)} \arrow[r,"d'"] \arrow[d,"\eta_{S^{(I)}}"'] &
S^{(J)} \arrow[r,"q'"] \arrow[d,"\eta_{S^{(J)}}"'] &
N \arrow[r] \arrow[d,"\eta_N"'] & 0\\
FG(S^{(I)}) \arrow[r,"FG(d')"'] &
FG(S^{(J)}) \arrow[r,"FG(q')"'] &
FG(N) \arrow[r] & 0
\end{tikzcd}
\]
前两列的竖直映射为同构，第三列是它们在余核上诱导的映射，故 \(\eta_N\) 也为同构。

所以 \(\eta:\operatorname{Id}\rightarrow FG\)、\(\varepsilon:GF\rightarrow\operatorname{Id}\) 都是自然同构，正好给出两函子互为拟逆。任意等价的反方向已经由\cref{b2:thm:morita-converse}证明，得到最后的充要条件。
\end{proof}

这就是\term[Morita-dengjia-dingli]{Morita 等价定理}[Morita equivalence theorem]。仅在正则模或有限自由模上检验单位、余单位，不足以替代证明中通过任意直和和余核所作的扩展。

\ProseParagraphBreak
同态也可显式恢复。若 \(a:F(M)\rightarrow F(M')\) 左 \(S\) 线性，则它唯一来自
\[
\varepsilon_{M'}\,G(a)\,\varepsilon_M^{-1}:M\longrightarrow M'.
\]
确实，式\eqref{b2:eq:morita-triangles}给出 \(F(\varepsilon_M)=\eta_{F(M)}^{-1}\)，单位的自然性便说明将上式施加 \(F\) 后得到 \(a\)。唯一性来自等价的忠实性。这使对象对应与同态对应同时可计算。

\begin{corollary}[互逆双模]\label{b2:cor:morita-inverse-bimodules}
设 \(R\) 为含幺环，\({}_RP\) 为幺左模范畴的有限生成投射生成元，\(S=\operatorname{End}_R(P)^{\mathrm{op}}\)，并令 \(ps=\alpha_s(p)\)，其中 \(\alpha_s\in\operatorname{End}_R(P)\) 对应 \(s\in S\)。令 \(Q=P^*=\operatorname{Hom}_R(P,R)\)，其幺 \((S,R)\) 双模作用为 \((s\lambda)(p)=\lambda(ps)\)、\((\lambda r)(p)=\lambda(p)r\)。则有双模同构
\[
{}_RP\otimes_S Q_R\cong {}_RR_R,\qquad
{}_SQ\otimes_RP_S\cong {}_SS_S.
\]
第一式把 \(p\otimes\lambda\) 送到 \(\lambda(p)\)；第二式把 \(\lambda\otimes p\) 送到 \(S\) 中对应于 \(x\mapsto\lambda(x)p\) 的元素。
\end{corollary}
\begin{proof}
第一式是\eqref{b2:eq:morita-counit}在 \(M=R\) 时的余单位，由\cref{b2:thm:morita-equivalence}为左 \(R\) 模同构。它也保持右 \(R\) 作用，因为 \((\lambda r)(p)=\lambda(p)r\)。

由\cref{b2:prop:finite-projective-hom-tensor}取 \(M=P\)，映射
\[
Q\otimes_RP\longrightarrow\operatorname{End}_R(P),\qquad
\lambda\otimes p\longmapsto T_{\lambda,p},\quad
T_{\lambda,p}(x)=\lambda(x)p
\]
为同构。有限对偶基解释了为什么这些自同态足以构成所有自同态。由\cref{b2:prop:finite-dual-basis}取 \((p_i,\lambda_i)\)，则 \(x=\sum_i\lambda_i(x)p_i\)，因而每个 \(R\) 线性自同态 \(T\) 都满足
\[
T(x)=\sum_i\lambda_i(x)T(p_i).
\]
所以任意自同态都是上述形式的有限和，且已证同构的逆将 \(T\) 送到 \(\sum_i\lambda_i\otimes T(p_i)\)。

将自同态 \(T\) 看作反环 \(S\) 的元素时，左、右 \(S\) 作用分别对应于两种复合次序：
\[
T_{s\lambda,p}=T_{\lambda,p}\circ\alpha_s,\qquad
T_{\lambda,ps}=\alpha_s\circ T_{\lambda,p}.
\]
由于 \(S\) 的乘法与自同态复合反向，这两式分别是 \(S\) 中的左乘 \(s\) 与右乘 \(s\)。故第二式是 \((S,S)\) 双模同构；这里 \(\operatorname{End}_R(P)\) 与 \(S\) 的识别用的是其加性群及这两个作用，而不是把两者的环乘法认作相同。
\end{proof}

这两个同构说明 \({}_RP_S\) 与 \({}_SQ_R\) 在双模张量意义下互为逆：它们的两种复合分别恢复正则双模 \({}_RR_R\) 与 \({}_SS_S\)。同一对双模因此也能从右模一侧构造等价。

\begin{corollary}\label{b2:cor:morita-left-right}
对任意含幺环 \(R,S\)，全部幺左模范畴等价当且仅当全部幺右模范畴等价：
\[
R\text{-}\mathbf{Mod}\simeq S\text{-}\mathbf{Mod}
\quad\Longleftrightarrow\quad
\mathbf{Mod}\text{-}R\simeq\mathbf{Mod}\text{-}S.
\]
\end{corollary}
\begin{proof}
若左模范畴等价，由\cref{b2:thm:morita-converse,b2:cor:morita-inverse-bimodules}取互逆双模 \({}_RP_S\) 与 \({}_SQ_R\)。函子 \(-\otimes_RP\) 把幺右 \(R\) 模送到幺右 \(S\) 模，\(-\otimes_SQ\) 把幺右 \(S\) 模送到幺右 \(R\) 模。两次复合经张量结合律分别同构于 \(-\otimes_R(P\otimes_SQ)\cong-\otimes_RR\) 和 \(-\otimes_S(Q\otimes_RP)\cong-\otimes_SS\)，故互为拟逆。反过来，把右 \(R,S\) 模分别视为左 \(R^{\mathrm{op}},S^{\mathrm{op}}\) 模，对刚证的方向再用一次，便得到左模范畴的等价。
\end{proof}

\subsection{满幂等元与角环的 Morita 等价}
\label{b2:subsec:1603:03}

\begin{proposition}[满幂等角环]\label{b2:prop:full-idempotent-morita}
设 \(R\) 为含幺环，\(e\in R\) 满足 \(e^2=e\)，令 \(S=eRe\)，其单位元为 \(e\)。条件 \(ReR=R\)、幺左 \(R\) 模 \(Re\) 为生成元，以及对每个幺左 \(R\) 模 \(M\) 都有 \(eM=0\Longrightarrow M=0\)，三者等价。在这些条件下，\(Re\) 为有限生成投射生成元，且
\[
R\text{-}\mathbf{Mod}\simeq S\text{-}\mathbf{Mod}
\]
由 \(M\mapsto eM\) 与 \(N\mapsto Re\otimes_SN\) 实现，其中 \(eM\) 取自然的幺左 \(S\) 模结构。对每个幺左 \(R\) 模 \(M\)，求值映射
\[
\varepsilon_M:Re\otimes_SeM\longrightarrow M,\qquad
re\otimes m\longmapsto rm
\]
为自然的左 \(R\) 模同构。
\end{proposition}
\begin{proof}
由 \(R=Re\oplus R(1-e)\)，\(Re\) 是正则左模的直和因子，故有限生成且投射。任意左模同态 \(Re\rightarrow R\) 由 \(e\) 的像 \(b\in eR\) 决定，公式为 \(re\mapsto rb\)。因此若 \(Re\) 为生成元，就有某个满射 \((Re)^{(L)}\rightarrow R\)。给 \(1\) 取一个原像，这个原像仅有有限支集；结合各分量上的同态公式便得到 \(1=\sum_{i=1}^n a_i e b_i\)，所以 \(ReR=R\)。此处和式有限来自这个原像的有限支集。

反过来，若 \(1=\sum_{i=1}^n a_i e b_i\)，映射 \((Re)^n\rightarrow R\)、\((p_i)\mapsto\sum_i p_i b_i\) 的像包含一，因而满射。其截面可取 \(r\mapsto(ra_i e)_i\)，所以 \(R\) 是 \((Re)^n\) 的直和因子。\cref{b2:thm:generator-characterizations}使 \(Re\) 为生成元。

若 \(ReR=R\) 且 \(eM=0\)，对每个 \(m\in M\) 有 \(m=\sum_i a_i e b_i m=0\)，所以 \(M=0\)。若 \(ReR\ne R\)，则 \(M=R/ReR\) 为非零幺左模，而 \(eM=0\)。这就证明第三个条件与满性等价，也说明非满的取角会把某个非零模送到零。

由\cref{b2:prop:corner-endomorphism}，\(\operatorname{End}_R(Re)^{\mathrm{op}}\cong S\)。与该命题中的同态识别一样，一个 \(f:Re\rightarrow M\) 完全由 \(f(e)\) 决定，因为 \(f(re)=rf(e)\)，且 \(f(e)=ef(e)\in eM\)。因此有自然同构
\[
\operatorname{Hom}_R(Re,M)\longrightarrow eM,\qquad f\longmapsto f(e),
\]
逆为 \(m\mapsto(re\mapsto rm)\)。若 \(re=r'e\)，因 \(m=em\) 有 \(rm=r'm\)，故逆公式良定义。对 \(s\in S\)，右乘 \(s\) 是 \(Re\) 的自同态，而 Hom 中的左作用满足 \((s\cdot f)(e)=f(es)=sf(e)\)，所以此同构保留左 \(S\) 作用。这里 \(eM\) 一般不是 \(M\) 的左 \(R\) 子模，所用的标量环是 \(S\)。当 \(e\) 满时，\(Re\) 已是有限生成投射生成元，\cref{b2:thm:morita-equivalence}给出题述等价。

求值映射 \(\varepsilon_M\) 也有直接的逆公式。固定一次表达 \(1=\sum_i a_i e b_i\)，则 \(m=\sum_i a_i e b_i m\)；要使求值恢复 \(m\)，就把每一项拆成 \(a_i e\otimes e b_i m\)。由此得到加性映射
\begin{equation}\label{b2:eq:full-idempotent-evaluation-inverse}
\delta_M(m)=\sum_i a_i e\otimes e b_i m.
\end{equation}
于是 \(\varepsilon_M\delta_M(m)=\sum_i a_i e b_i m=m\)。反向复合在 \(re\otimes m\)、\(m=em\) 上得到
\[
\sum_i a_i e\otimes e b_i r e m
=\sum_i a_i e b_i r e\otimes m=re\otimes m,
\]
其中中间一步移动的平衡标量是 \(e b_i r e\in S\)。故 \(\delta_M\) 与左 \(R\) 线性映射 \(\varepsilon_M\) 互为逆，\(\delta_M\) 也左 \(R\) 线性。求值与模同态相容，因而此同构自然。另一方面，\(e(Re\otimes_SN)\) 自然识别为 \(S\otimes_SN\cong N\)：右 \(S\) 模分裂投影 \(Re\rightarrow S\)、\(p\mapsto ep\) 与其包含映射张量后仍互为投影与截面，故它的像正是 \(e(Re\otimes_SN)\)，最后使用单位张量同构。
\end{proof}

对含幺环 \(B\) 及 \(n\ge1\)，在矩阵环 \(R=M_n(B)\) 中，\(e=E_{11}\) 满，因为 \(1=\sum_iE_{i1}eE_{1i}\)。上述公式正好恢复\cref{b2:prop:matrix-morita-equivalence}。当 \(B=k\) 为域且 \(n>1\) 时，令 \(A=M_n(k)\)，则幺左 \(A\) 模 \(Ae\) 是有限生成投射生成元，却不是自由模：\(\dim_kAe=n\)，而任意非零自由左 \(A\) 模都包含一个正则模 \(A\) 的直和分量，故其 \(k\) 维数至少为 \(n^2>n\)，无论自由基有限还是无限都不可能同构于 \(Ae\)。

\ProseParagraphBreak
若 \(R=B\times C\)、\(e=(1,0)\)，且 \(C\ne0\)，则 \(ReR=B\times0\ne R\)，取角只保留第一个坐标；第二个坐标上的非零模被送到零，所以不能给出等价。

\ProseParagraphBreak
对有限维代数 \(A\)，\cref{b2:prop:basic-corner}从每种有限生成不可分解投射模的同构类型中选取一份，构造了满足 \(AeA=A\) 的幂等元，并证明 \(eAe\) 为基本代数。\cref{b2:prop:full-idempotent-morita}进一步说明 \(A\) 与这个基本代数的全部幺左模范畴等价。因此可先在基本代数上研究模，再经上述张量函子返回原代数。

\begin{proposition}[满角环的双边理想对应]\label{b2:prop:full-corner-ideal-correspondence}
设 \(R\) 为含幺环，\(e\in R\) 为满幂等元，即 \(e^2=e\) 且 \(ReR=R\)，令 \(S=eRe\)。对双边理想 \(I\subseteq R\) 与 \(K\subseteq S\)，映射
\[
I\longmapsto eIe,\qquad K\longmapsto RKR
\]
给出 \(R\) 与 \(S\) 的双边理想格之间互逆且保持包含关系的对应，其中 \(RKR\) 表示元素 \(rkt\)、\(r,t\in R\)、\(k\in K\) 的有限和组成的双边理想。对 \(R\) 的每个双边理想 \(I\)，记 \(\bar e=e+I\)，则 \(\bar e\) 是 \(R/I\) 的满幂等元，并有环同构
\[
\bar e(R/I)\bar e\cong S/eIe.
\]
因此 \(R/I\) 与 \(S/eIe\) 的全部幺左模范畴等价。
\end{proposition}
\begin{proof}
若 \(I\) 为 \(R\) 的双边理想，则 \(eIe\) 为 \(S\) 的双边理想；若 \(K\) 为 \(S\) 的双边理想，则上述有限和组成的 \(RKR\) 为 \(R\) 的双边理想。固定 \(1=\sum_{i=1}^n a_i e b_i\)。对 \(x\in I\)，两次使用这个等式得到
\[
x=\sum_{i,j}a_i e(b_i x a_j)e b_j\in R(eIe)R.
\]
反向包含 \(R(eIe)R\subseteq I\) 来自 \(I\) 的双边理想性质，故 \(R(eIe)R=I\)。

对 \(k\in K\) 有 \(k=eke\)，所以 \(K\subseteq e(RKR)e\)。反向包含可逐项检查：若 \(r,t\in R\)、\(k\in K\)，则
\[
e(rkt)e=(ere)k(ete)\in K.
\]
故 \(e(RKR)e=K\)。两种映射显然都保持包含关系，因此给出互逆的理想格对应。

商映射限制到 \(S\) 给出满环同态
\[
S\longrightarrow\bar e(R/I)\bar e,\qquad
ere\longmapsto\bar e(r+I)\bar e.
\]
其核为 \(S\cap I=eIe\)，因为 \(s\in S\cap I\) 时 \(s=ese\)。商环同构定理便给出所述角环同构。\(1=\sum_i a_i e b_i\) 在 \(R/I\) 中仍成立，故 \(\bar e\) 满；由\cref{b2:prop:full-idempotent-morita}应用于 \(R/I\) 即得模范畴等价。此论证也包含 \(I=R\) 的情形，此时两个商环均为零环。
\end{proof}

% !TeX root = ../../Book2.tex
% 纲要标题：### 16.4 Morita 不变量
% 纲要位置：计划纲要.md，第 1583 行。

\section{Morita 不变量}
\label{b2:sec:1604}

模范畴的等价保留对象、同态及其相互关系，却不把一个模的底集合或坐标照搬过去。若一种性质还使用了指定的正则模、到向量空间的忘却函子或额外的型，就须另行检查这些数据是否相容。例如，相应模的合成长度数值确实保持，而它们作为同一个底域上的向量空间，维数可以改变；两个正则模也未必彼此对应。中心和子模结构则可以从模范畴本身恢复。

% 纲要要点（供目录核对）：
%
% * 中心、简单模、投射模和内射模；中心作为恒等函子自然自变换环，子模作为单射子对象及其因子关系
% * 半单性、不可分解性与有限长度；正式证明有限展示与环的左 Noether、左 Artin 性保持，以及与除环等价的单 Artin 判据
% * 维数、生成元数及附加型不是 Morita 不变量
% * 在同一实向量空间上比较正定与不定型及其不同伴随运算，说明普通 Morita 数据没有指定型
% ---
%

\subsection{中心、简单模与投射性}
\label{b2:subsec:1604:01}

\begin{proposition}\label{b2:prop:morita-center}
设 \(R,S\) 为含幺环，幺性左模范畴 \(\mathcal C_R=R\text{-}\mathbf{Mod}\) 与 \(\mathcal C_S=S\text{-}\mathbf{Mod}\) 等价。对 \(T=R,S\)，恒等函子的自然自变换以逐分量加法、复合作为乘法、恒等自然变换作为单位，组成环
\(\operatorname{Nat}(\operatorname{Id}_{\mathcal C_T},\operatorname{Id}_{\mathcal C_T})\)。映射
\[
Z(T)\longrightarrow
\operatorname{Nat}(\operatorname{Id}_{\mathcal C_T},\operatorname{Id}_{\mathcal C_T}),
\qquad
c\longmapsto\bigl(\theta^c_M:m\longmapsto cm\bigr)_{M\in\mathcal C_T}
\]
是含幺环同构。因此 \(Z(R)\cong Z(S)\)。
\end{proposition}
\begin{proof}
中心与自然自变换环的同构已由\cref{b2:prop:natural-center}证明。其逆映射是在正则模的一处求值：\(c=\theta_T(1)\)。任意模元素 \(m\in M\) 都给出同态 \(f_m:T\to M\)、\(r\mapsto rm\)，自然性强迫 \(\theta_M(m)=cm\)，因此这一求值决定所有分量；加法、复合与单位分别对应 \(c+d\)、\(cd\) 与 \(1\)。

接下来将一个范畴中对所有对象自然定义的自同态，传到另一个范畴。设 \(H:\mathcal C_R\rightarrow\mathcal C_S\) 为等价，固定拟逆 \(K\) 及自然同构 \(\beta:HK\Rightarrow\operatorname{Id}_{\mathcal C_S}\)。给定 \(\theta\in\operatorname{Nat}(\operatorname{Id}_{\mathcal C_R},\operatorname{Id}_{\mathcal C_R})\)，对每个幺性左 \(S\) 模 \(Y\) 令
\[
\theta'_Y=\beta_Y\,H(\theta_{K(Y)})\,\beta_Y^{-1}.
\]
\(\beta\) 与 \(\theta\) 的自然性保证这些分量构成自然自变换。由\cref{b2:lem:module-equivalence-structure}，\(H\) 加性，故此对应保持逐分量加法；共轭保持复合与单位。对 \(H(X)\)，全忠实性使 \(\beta_{H(X)}=H(b_X)\)，其中 \(b_X:KH(X)\to X\) 为同构。\(\theta\) 对 \(b_X\) 的自然性给出
\[
\theta'_{H(X)}
=H\bigl(b_X\theta_{KH(X)}b_X^{-1}\bigr)
=H(\theta_X).
\]
因此这个对应单射。反过来，目标范畴上的任一自然自变换 \(\tau\) 在各 \(H(X)\) 上的分量可由全忠实性唯一提升为 \(\theta_X\)。\(\tau\) 对 \(H(f)\) 的自然性与 \(H\) 的忠实性保证 \(\theta\) 自然。它的像与 \(\tau\) 在所有 \(H(X)\) 上相同，再用同构 \(\beta_Y:HK(Y)\to Y\) 的自然性，就知在每个 \(Y\) 上也相同。因此两个自然自变换环同构，进而 \(Z(R)\cong Z(S)\)。
\end{proof}

\begin{proposition}\label{b2:prop:morita-submodule-lattices}
设 \(R,S\) 为含幺环，\(H:R\text{-}\mathbf{Mod}\rightarrow S\text{-}\mathbf{Mod}\) 为幺性左模范畴的等价，\(M\) 为幺性左 \(R\) 模，则给定的 \(H\) 诱导 \(M\) 与 \(H(M)\) 的子对象格之间的序同构。把子对象视为子模时，\(N\subseteq M\) 对应于 \(H(N\hookrightarrow M)\) 的像；这一序同构保留零子模、全模、任意子模族的交与和。
\end{proposition}
\begin{proof}
子对象由单射 \(i:N\to M\) 表示，两个单射在定义域之间存在与它们相容的同构时，表示同一子对象。在模范畴中，取 \(i(N)\subseteq M\) 就给出它的子模代表。\cref{b2:lem:module-equivalence-structure}保证 \(H(i)\) 仍单射，所以其像同构于 \(H(N)\)；这里对应的是嵌入的像，而不是对底集合逐元素应用 \(H\)。反过来，给定 \(L\subseteq H(M)\)，选同构 \(u:H(X)\rightarrow L\)，把复合 \(H(X)\rightarrow L\hookrightarrow H(M)\) 唯一提升为 \(a:X\rightarrow M\)。等价反映单射，故 \(a\) 单射，令 \(N=a(X)\)，便得到所需原子模。若两个原子模给出相同像，像之间保持嵌入的同构提升后仍与原嵌入相容，因而两者在 \(M\) 中的像相同，故子模相等。

子模包含等价于两个嵌入之间存在相容的因子映射。全忠实性双向保持这一条件，所以所得双射保持并反映包含关系。零与全模是最小、最大元；交为最大下界，子模之和为最小上界，因此它们也由这个序同构保留。
\end{proof}

非零模简单，当且仅当子模序只有零和全模两个元素，所以简单模在等价下对应。\cref{b2:lem:module-equivalence-structure}已经从满射提升证明投射性也保持。并不要求简单模对应同一个底集合，或投射模仍表现为相同秩的自由模。

\subsection{半单性与有限长度}
\label{b2:subsec:1604:02}

\begin{proposition}\label{b2:prop:morita-invariants}
设 \(R,S\) 为含幺环，\(H:R\text{-}\mathbf{Mod}\rightarrow S\text{-}\mathbf{Mod}\) 为幺性左模范畴的等价，则 \(H\) 保持并反映模的简单性、投射性、内射性、不可分解性、半单性、有限生成性、有限展示性、Noether 性、Artin 性和有限长度。若幺性左 \(R\) 模 \(M\) 有限长度，则
\[
\ell_R(M)=\ell_S\bigl(H(M)\bigr).
\]
相应简单因子的合成重数也保持。此外，\(R\) 为左 Noether 环，当且仅当 \(S\) 为左 Noether 环；左 Artin 性也满足同样的等价。两个环还同时为半单环或同时不是半单环。
\end{proposition}
\begin{proof}
简单性和投射性已经证明，内射性由\cref{b2:lem:module-equivalence-structure}得到。一个模分解为两个非零模的直和，等价于存在相应的嵌入、投影及直和恒等式；这也是\cref{b2:prop:idempotent-direct-sum}中非平凡幂等自同态的投影刻画。等价双向保持这些映射关系和非零对象，故保持并反映不可分解性。半单模是简单模的某个直和；等价保持简单模与任意直和，所以保持半单性，对拟逆使用同样结论得到反映性。

由\cref{b2:prop:morita-submodule-lattices}，子模链的包含、严格性和稳定性双向保持，因此 Noether 性与 Artin 性保持。对一条合成列，等价的正合性使相邻商对应原简单因子的像；每个像仍简单，故得到相同层数的合成列。\cref{b2:thm:module-jordan-holder}使这个层数正是两边的长度，不同简单类型也由全忠实性一一对应，所以重数相同。对拟逆应用同样论证，有限长度也被反映。

有限生成性不能通过比较两边的生成元个数来证明，但可以从子模序直接判定：一个模 \(M\) 有限生成，当且仅当每个非空、向上有向且并为 \(M\) 的子模族中，都有一个成员等于 \(M\)。向上有向是指任意两个成员都包含在某个共同成员中。正向时，每个有限生成元属于族中某个成员，反复使用有向性便找到包含全部生成元的共同成员，它只能等于 \(M\)；若 \(M=0\)，任取一个成员即可。反向取 \(M\) 的所有有限生成子模，它们向上有向，因为两组有限生成元合在一起仍有限，且并为 \(M\)，因为每个元素都属于自己的循环子模。于是其中一个成员等于 \(M\)。子模序同构保持有向性和最小上界；对有向子模族，这个最小上界正是它们的并，因此保持这一判据。

有限展示还须检查全部关系能否有限生成。令 \(P=H({}_RR)\)，已证的有限生成性与投射性的保持使 \(P\) 为有限生成投射左 \(S\) 模。若 \(M\) 有有限自由展示 \(R^a\xrightarrow{d}R^b\twoheadrightarrow M\to0\)，其中 \(a,b\ge0\) 为整数，正合性与有限直和的保持给出
\[
P^a\xrightarrow{H(d)}P^b\twoheadrightarrow H(M)\longrightarrow0.
\]
由\cref{b2:prop:finite-projective-summand}，可取有限生成投射左 \(S\) 模 \(Q\)，使 \(P^b\oplus Q\cong S^t\)，其中 \(t\ge0\) 为整数。把 \(P^b\) 上的满射与 \(Q\) 上的零映射合起来，得到有限自由满射 \(S^t\to H(M)\)，其核在这个分解下为
\[
\operatorname{im}H(d)\oplus Q.
\]
第一项是有限生成模 \(P^a\) 的像，第二项也有限生成，故这个核有限生成。为核取一个有限自由满射，就得到 \(H(M)\) 的有限自由展示。对拟逆应用同样论证，有限展示性也被反映。

环的链条件还需要从 \(P\) 传到正则模 \({}_SS\)，因为两者未必同构。由\cref{b2:lem:module-equivalence-structure}，\(P\) 也是生成元，因而\cref{b2:prop:finite-projective-summand,b2:thm:generator-characterizations}分别给出有限直和分解
\[
P\oplus Q_1\cong S^u,\qquad
S\oplus Q_2\cong P^v,
\]
其中 \(u,v\ge0\) 为整数，\(Q_1,Q_2\) 为左 \(S\) 模。\cref{b2:thm:chain-conditions-exact}保证 Noether 性和 Artin 性对有限直和及子模传递。第一个分解使 \({}_SS\) 的链条件传到 \(P\)，第二个分解使 \(P\) 的链条件传到 \({}_SS\)。而子模格对应已说明 \(P=H(R)\) 与 \({}_RR\) 的链条件相同，所以两环同时为左 Noether 环，也同时为左 Artin 环。

最后，\cref{b2:thm:semisimple-ring-characterizations}说明环半单当且仅当它的每个左模半单。每个目标对象都同构于等价像中的对象，而半单性又被双向保持，所以两环的这一性质同时成立。这里不能只比较正则模的像，因为等价未必把一个正则模送到另一个正则模。
\end{proof}

有限长度并不要求两边的环都有限维；它由子模链定义，正是等价能保留的结构。例如，若 \(M\) 的两个不可分解直和分量合成长度分别为三和五，它们在等价后仍不可分解，长度仍分别为三和五：分解与非零性都可由投影和零对象辨认，不能因坐标数改变而改变这些长度。

\ProseParagraphBreak
中心同构也给出一个实用限制：两个交换环若 Morita 等价，则它们作为环同构，因为交换环就是自己的中心。反向，环同构当然给出改写标量的模范畴同构。因此对交换环，Morita 等价不会把不同的环合并；其更广的作用主要体现在非交换结构中。

\subsection{维数、生成元数与附加型}
\label{b2:subsec:1604:03}

设 \(k\) 为域、\(n>1\) 为整数，取 \(R=k\)、\(S=M_n(k)\)。它们 Morita 等价，但作为 \(k\) 代数的维数分别为一和 \(n^2\)，前者交换，后者不交换。标准等价把 \(k\) 模 \(V\) 送到列模 \(V^n\)，在同态上逐坐标作用，是 \(k\) 线性的。即使如此，若 \(\dim_kV=m<\infty\)，仍有
\[
\dim_kV=m,\qquad \dim_k(V^n)=nm.
\]
维数变化并没有违反长度不变性：\(V\) 是 \(m\) 个一维简单 \(k\) 模的直和，\(V^n\) 是 \(m\) 个简单 \(M_n(k)\) 列模的直和，故两边的模长度都为 \(m\)。

\ProseParagraphBreak
最少生成元的数值也可能改变。对 \(V=k^m\)，有 \(\mu_k(V)=m\)；在矩阵环一侧，\cref{b2:prop:semisimple-top-generators}给出
\[
\mu_{M_n(k)}(V^n)=\left\lceil\frac mn\right\rceil.
\]
因此等价保持“有限生成”这个性质，却不保持相对于正则模计算的最少生成元数。事实上，
\[
{}_{M_n(k)}M_n(k)\cong(k^n)^{\oplus n},
\]
所以一个正则模包含 \(n\) 个简单列模。一个生成元给出的正则模像至多覆盖 \(n\) 份简单列模；要覆盖 \(V^n\) 中的 \(m\) 份，最少需要 \(\lceil m/n\rceil\) 个生成元。

\ProseParagraphBreak
若 \(k\) 为 \(q\) 元域，则两个环的元素数分别为 \(q\) 和 \(q^{n^2}\)。模的元素从向量换成向量列，环元素从标量换成矩阵，都允许底集合发生变化。使用 Morita 等价时，应说明所比较的是环同构、模范畴结构、模长度还是指定底域上的维数；这些问题不能仅凭“等价”二字互相替换。

\begin{example}\label{b2:exm:morita-forgets-forms}
普通模范畴不记录一个模上额外指定的型。取 \(V=\mathbb R^2\)，在同一个实向量空间上分别指定矩阵为
\[
H_+=I_2,\qquad H_-=\operatorname{diag}(1,-1)
\]
的对称双线性型。它们的底层模完全相同，自同态代数也都是 \(M_2(\mathbb R)\)，但两型不等距：第二个型有非零各向同性向量 \((1,1)\)，第一个没有。
\end{example}

\cref{b2:prop:form-adjoint-operator}将这些型分别变为自同态代数上的伴随运算
\[
\sigma_+(X)=X^{\mathsf T},\qquad
\sigma_-(X)=H_-^{-1}X^{\mathsf T}H_-.
\]
它们对 \(E_{12}\) 的取值分别为 \(E_{21}\) 与 \(-E_{21}\)，故是不同运算。Morita 等价中的对象和态射没有包括 \(H_+\)、\(H_-\) 或这些伴随运算，因而不能仅凭模范畴等价，就声称还给出了带型对象之间的等价。若要比较后者，必须额外指定如何传送型、对偶及等距映射；普通 Morita 定理并未给出这些附加数据。

\ProseParagraphBreak
回到半单环，Artin–Wedderburn 分解中的矩阵大小可以在模范畴等价下改变，而简单模的同构类型仍一一对应。这给出了与除环等价的准确范围。

\begin{corollary}\label{b2:cor:morita-division-ring}
设 \(R\) 为非零含幺环。\(R\) Morita 等价于某个除环，当且仅当 \(R\) 是单 Artin 环，即 \(R\) 是单环且为左 Artin 环。
\end{corollary}
\begin{proof}
若 \(R\) 为单 Artin 环，\cref{b2:prop:simple-artinian}给出 \(R\cong M_n(D)\)，其中 \(D\) 为除环、\(n\ge1\)。由\cref{b2:prop:matrix-morita-equivalence}，这个矩阵环与 \(D\) Morita 等价。

反过来，设 \(R\text{-}\mathbf{Mod}\simeq D\text{-}\mathbf{Mod}\)，其中 \(D\) 为除环。每个幺性左 \(D\) 模都有基，因而是简单模 \({}_DD\) 的直和，且非零简单左 \(D\) 模只有这一种同构类型。\cref{b2:prop:morita-invariants}使每个幺性左 \(R\) 模也半单，并使简单左 \(R\) 模只有一种同构类型。由\cref{b2:thm:semisimple-ring-characterizations}，\(R\) 为半单环。再用\cref{b2:thm:artin-wedderburn}，将 \(R\) 写成有限个除环上矩阵环的乘积。不同矩阵因子的简单列模不同构，因为相应的中心幂等元在一个上作用为恒等，在另一个上作用为零，所以这里只有一个因子。由于 \(R\ne0\)，这个因子存在，由\cref{b2:prop:simple-artinian}即得 \(R\) 为单 Artin 环。
\end{proof}

只有一种简单模类型并不足够。设 \(k\) 为域，\(A=k[\varepsilon]/(\varepsilon^2)\)。对任意简单左 \(A\) 模 \(U\)，\(\varepsilon U\) 是子模；若它等于 \(U\)，则 \(U=\varepsilon U=\varepsilon^2U=0\)，矛盾，故 \(\varepsilon U=0\)。于是 \(U\) 是 \(k\) 上的简单模，同构于 \(k\)。但正则模 \(A\) 并不半单：商映射 \(A\to A/(\varepsilon)\) 若有模截面，一的像应为某个 \(1+c\varepsilon\)，并被 \(\varepsilon\) 零化；实际上 \(\varepsilon(1+c\varepsilon)=\varepsilon\ne0\)。所以这个商映射不能分裂，\(A\) 不与除环 Morita 等价。


\begin{chaptersummary}
生成元使每个模都成为其某个直和的商，有限生成使 Hom 保持任意直和，投射性使 Hom 正合。三者合起来，允许从正则对象上的伴随同构出发，经由展示和余核覆盖全部模。反方向，等价将正则模送到投射生成元，再由紧性恢复有限生成，所以这一构造恰好描述全部 Morita 等价。

\ProseParagraphBreak
自同态环的反环记录左模与右作用之间的乘法次序。互逆双模还给出相应的右模范畴等价。矩阵环等价和满幂等角环等价都能写成明确的矩阵单位及求值公式；满性保证没有一部分模在取角时被丢弃。基本角代数因而保留原有限维代数的整个模范畴。

\ProseParagraphBreak
中心可由恒等函子的自然自变换恢复，子模序则由单态射及其因子关系恢复。这些描述解释了简单性、投射性、内射性、不可分解性、半单性、有限生成性、链条件和长度为什么保持。有限展示性还须检查关系模的有限生成，环的链条件则须把正则模的像与另一侧正则模联系起来；这些结论已在正文中分别证明。与此同时，正则模未必对应正则模，底域维数、元素数和最少生成元数都可能改变。额外指定的型及伴随运算也不是普通模范畴所记录的结构，若要传送这些数据，必须另加相容条件。后续比较中心单代数时，将继续使用这种区分，而不把 Morita 等价误当作环同构。
\end{chaptersummary}

\BookExercises

\begin{exercise}\label{b2:ex:16-trace-generator}
设 \(k\) 为域、\(R=k[t]/(t^3)\)，以 \(t\) 表示剩余类。分别计算 \(R/(t)\)、\(R/(t^2)\) 及 \(P=R/(t)\oplus R/(t^2)\) 的迹理想，判断 \(P\) 是否为生成元。再判断 \(P\oplus R\) 是否为生成元。
\end{exercise}
\begin{solution}
对任意理想 \((a)\)，从 \(R/(a)\) 到 \(R\) 的同态由 \(1+(a)\) 的像 \(b\) 决定，条件恰为 \(ab=0\)。因而在当前交换环中，迹理想为 \(\operatorname{Ann}_R(a)\)。这里
\[
\operatorname{tr}_R(R/(t))=(t^2),\qquad
\operatorname{tr}_R(R/(t^2))=(t).
\]
每个从有限直和出发的同态由各分量同态相加，所以 \(\operatorname{tr}_R(P)=(t)\ne R\)。由\cref{b2:prop:trace-generator-criterion}，\(P\) 不是生成元。另一方面，\(P\oplus R\) 含正则模为直和因子，其迹包含 \(\operatorname{id}_R(R)=R\)，所以是生成元。
\end{solution}

\begin{exercise}\label{b2:ex:16-object-detection}
设 \(k\) 为域、\(n\ge2\)，令 \(R=k[t]/(t^n)\)、\(P=R/(t)\)。对任意左 \(R\) 模 \(M\)，识别 \(\operatorname{Hom}_R(P,M)\)，并求从全部 \(P\) 副本到 \(M\) 的求值映射之像。由此说明 \(P\) 能检测非零对象却不能检测所有非零同态，并刻画哪些 \(M\) 是 \(P\) 某个直和的商。
\end{exercise}
\begin{solution}
在 \(1+(t)\) 处求值得到
\[
\operatorname{Hom}_R(P,M)\cong\{m\in M:tm=0\}.
\]
这也是所有求值像的和，因为每个被 \(t\) 消去的元素都给出一个同态。若 \(M\ne0\)，取 \(x\ne0\)，并取最小正整数 \(j\) 使 \(t^jx=0\)；它存在，因为 \(t^nM=0\)。则 \(t^{j-1}x\ne0\) 被 \(t\) 消去，故 Hom 非零。然而 \(R\) 上乘 \(t\) 是非零同态，却在所有来自 \(P\) 的像上为零。这是\cref{b2:prop:object-detection-not-generator}在较高幂零次数中的同一现象。

\(P\) 的任意直和都被 \(t\) 消去，其商也如此。反过来，若 \(tM=0\)，则 \(M\) 是一个 \(k\) 向量空间；取一组基便得到若干 \(P\cong k\) 的直和与 \(M\) 的同构。因此题述商恰为被 \(t\) 消去的模。
\end{solution}

\begin{exercise}\label{b2:ex:16-projective-not-generator}
设 \(k\) 为域、\(R=k\times k\times k\)，令 \(P=k^2\oplus k\oplus0\)，三个因子分别在对应分量上作用。取第一分量的行坐标，将 \(\operatorname{End}_R(P)^{\mathrm{op}}\) 识别为 \(S=M_2(k)\times k\)。对 \(M=M_1\oplus M_2\oplus M_3\)，写出 \(\operatorname{Hom}_R(P,M)\) 的两个 \(S\) 分量，并求余单位 \(P\otimes_S\operatorname{Hom}_R(P,M)\to M\) 的像。判断哪些模被 Hom 函子送到零。
\end{exercise}
\begin{solution}
三个中心幂等元迫使同态分别作用于对应分量，故
\[
\operatorname{Hom}_R(P,M)\cong M_1^2\oplus M_2
\]
作为左 \(M_2(k)\times k\) 模，第一项为矩阵乘列，第二项为通常的标量作用。第一项的行列求值由\cref{b2:prop:matrix-row-column-bimodules}恢复 \(M_1\)，第二项恢复 \(M_2\)，所以余单位的像恰为 \(M_1\oplus M_2\oplus0\)。第三项不能由 \(P\) 产生；所有 \(0\oplus0\oplus M_3\) 都被送到零，且只有这些模被送到零。\(P\) 是有限生成投射模，但不是生成元，说明增加第一分量的副本并不能弥补缺失的第三分量。
\end{solution}

\begin{exercise}\label{b2:ex:16-generator-not-projective}
设 \(k\) 为域、\(R=k[t]/(t^2)\)、\(P=R\oplus R/(t)\)，令 \(q:R\to R/(t)\) 为商映射。证明 \(P\) 有限生成且为生成元，计算
\[
\operatorname{Hom}_R(P,R)\longrightarrow\operatorname{Hom}_R(P,R/(t))
\]
的像及两个 Hom 空间的 \(k\) 维数，并由此判断投射性。
\end{exercise}
\begin{solution}
\(P\) 有限生成，且含 \(R\) 为直和因子，所以是生成元。第一分量的同态由任意 \(r\in R\) 决定，第二分量的像必须被 \(t\) 消去，故
\[
\operatorname{Hom}_R(P,R)\cong R\oplus kt,\qquad
\operatorname{Hom}_R(P,R/(t))\cong k\oplus k.
\]
维数分别为三和二。后复合 \(q\) 把 \((r,ct)\) 送到 \((r+(t),0)\)，其像只有第一项，维数为一。特别地，从 \(P\) 到第二项 \(R/(t)\) 的投影不能沿 \(q\) 提升。因此 Hom 不保持这一满射，\(P\) 不投射。有限生成和生成元性质并不能替代提升条件。
\end{solution}

\begin{exercise}\label{b2:ex:16-projective-generator-not-compact}
设 \(k\) 为域、\(P=\bigoplus_{j\ge1}ke_j\)。用无限矩阵描述
\[
\bigoplus_{i\ge1}\operatorname{Hom}_k(P,k)
\longrightarrow\operatorname{Hom}_k\left(P,\bigoplus_{i\ge1}k\right)
\]
的像，说明“每列只有有限个非零条目”为什么不足以保证来自左边。
\end{exercise}
\begin{solution}
目标中的同态由各 \(e_j\) 的像确定，故对应每列只有有限个非零条目的矩阵 \((a_{ij})\)。左边是有限多个目标分量上的同态之和，所以对应的矩阵只有有限个非零行；每一行可以有无限多个非零条目。反过来，只有有限个非零行的矩阵确实给出左边的元素。

单位矩阵每列恰有一个非零条目，却有无限多个非零行，因而不在像中。这里有限支集可以随输入的 \(e_j\) 变化，紧性要求的是全像共有的有限目标支集。\(P\) 自由而且含一份 \(k\)，所以投射且为生成元；失败的正是有限生成性。
\end{solution}

\begin{exercise}\label{b2:ex:16-full-idempotent-detects}
设 \(k\) 为域、\(R=T_2(k)\)，分别取 \(e=E_{11}\)、\(f=E_{22}\)。计算 \(eRe\)、\(fRf\)、\(ReR\)、\(RfR\)，并求各自正则模上的求值映射之像。对每个角给出一个被送到零的非零简单模。
\end{exercise}
\begin{solution}
两个角都同构于 \(k\)，但
\[
ReR=kE_{11}\oplus kE_{12},\qquad
RfR=kE_{12}\oplus kE_{22},
\]
都是真理想。求值映射 \(Re\otimes_{eRe}eR\to R\) 的像是 \(ReR\)：纯张量的像为乘积 \(reb\)，其和恰为这个理想；另一角同理。这两张量源各有 \(k\) 维数二，两个矩阵单位的乘积分别给出所示理想的基，故这两个求值映射均单射，却不满射。

两个一维简单模 \(S_1,S_2\) 分别由第一、第二对角条目作用。\(eS_2=0\)，而 \(fS_1=0\)。所以即使角环与标量域相同、求值映射还单射，未满的角仍会遗漏非零对象。
\end{solution}

\begin{exercise}\label{b2:ex:16-rectangular-morita}
取 \(R=M_3(k)\)、\(e=E_{11}+E_{22}\)。证明 \(eRe\cong M_2(k)\)，给出满性的一次有限和表达，并写出对任意 \(M\) 的求值映射 \(Re\otimes_{eRe}eM\rightarrow M\) 的逆。在自然列模 \(k^3\) 上，这个等价得到哪个 \(M_2(k)\) 模？
\end{exercise}
\begin{solution}
角环是左上二阶矩阵块，单位为 \(e\)，所以同构于 \(M_2(k)\)。等式
\[
1=e+E_{31}eE_{13}
\]
证明 \(e\) 满。应用\cref{b2:prop:full-idempotent-morita}及其逆式\eqref{b2:eq:full-idempotent-evaluation-inverse}，代入 \((a_1,b_1)=(1,1)\)、\((a_2,b_2)=(E_{31},E_{13})\)，求值映射的逆为
\[
m\longmapsto e\otimes em+E_{31}\otimes E_{13}m.
\]
这里 \(E_{31}\in Re\)、\(E_{13}m\in eM\)；正向求值后得到 \(em+E_{33}m=m\)，反向复合由角环平衡关系核验。对 \(M=k^3\)，\(eM\) 为前两个坐标组成的自然列模 \(k^2\)。两个简单模的基域维数分别为三和二，但模长度都为一。
\end{solution}

\begin{exercise}\label{b2:ex:16-matrix-hom}
设 \(R\) 为含幺环、\(n\ge1\)，\(M,N\) 为幺左 \(R\) 模，令 \(A=M_n(R)\)。证明每个左 \(A\) 模同态 \(u:M^n\rightarrow N^n\) 都唯一地由一个左 \(R\) 模同态 \(h:M\rightarrow N\) 逐坐标作用得到。
\end{exercise}
\begin{solution}
把 \(m\in M\) 放到第一坐标，因 \(u\) 与 \(E_{11}\) 交换，其像只可能在第一坐标非零，记该坐标为 \(h(m)\)。与 \(rE_{11}\) 交换说明 \(h(rm)=r\cdot h(m)\)，加法相容也直接成立，所以 \(h\) 左线性。第 \(i\) 坐标向量由 \(E_{i1}\) 作用于第一坐标向量得到，因此
\[
u(0,\ldots,m,\ldots,0)=\bigl(0,\ldots,h(m),\ldots,0\bigr).
\]
有限加法得到 \(u(m_i)_i=\bigl(h(m_i)\bigr)_i\)。反向逐坐标的 \(h\) 保持矩阵作用，故确为 \(A\) 模同态；在第一坐标求值又保证唯一性。
\end{solution}

\begin{exercise}\label{b2:ex:16-row-column-evaluation}
设 \(B\) 为任意环，令 \(A=M_2(B)\)、\(S=M_3(B)\)，取矩形矩阵双模 \({}_AP_S=M_{2\times3}(B)\)、\({}_SQ_A=M_{3\times2}(B)\)。证明矩阵乘法给出双模同构
\[
P\otimes_SQ\cong A,\qquad Q\otimes_AP\cong S,
\]
并在矩阵单位上写出逆。
\end{exercise}
\begin{solution}
用 \(p_{i\ell}(b)\) 表示仅第 \((i,\ell)\) 项为 \(b\) 的 \(P\) 元素，\(q_{\ell j}(b)\) 类似。矩阵乘法结合律给出所需平衡性。第一映射的逆为
\[
(a_{ij})\longmapsto\sum_{i,j=1}^2p_{i1}(a_{ij})\otimes q_{1j}(1).
\]
确实，在 \(S\) 上移过 \(E_{1\ell}\)，再移过 \(cE_{11}\)，得到
\[
p_{i\ell}(b)\otimes q_{kj}(c)
=\delta_{\ell k}\,p_{i1}(bc)\otimes q_{1j}(1).
\]
每个纯张量可由这类矩阵单位张量相加，所以这个公式同时验证两个复合为恒等。第二映射的逆为
\[
(s_{\ell k})\longmapsto\sum_{\ell,k=1}^3q_{\ell1}(s_{\ell k})\otimes p_{1k}(1),
\]
在 \(A\) 上使用同样的平衡计算即可。所列映射保持外侧矩阵作用；系数始终按 \(bc\) 的次序相乘，不要求 \(B\) 交换。此题把\cref{b2:prop:matrix-row-column-bimodules}的行列模型推广到两个不同的矩阵大小。
\end{solution}

\begin{exercise}\label{b2:ex:16-dual-bimodule-pair}
设 \(R\) 为任意环、\(u\in R\)，在行列双模模型 \(P=R^{1\times2}\)、\(Q=R^{2\times1}\)、\(S=M_2(R)\) 中取
\[
C=\begin{pmatrix}1&u\\0&1\end{pmatrix}.
\]
若行坐标改为 \(\widetilde p=pC\)，求保持 \(pq\) 的列坐标变换及与右作用相容的矩阵坐标变换。对 \(B=\begin{psmallmatrix}a&b\\c&d\end{psmallmatrix}\) 写出新矩阵，核对两个双模配对。
\end{exercise}
\begin{solution}
因为 \(C^{-1}=\begin{psmallmatrix}1&-u\\0&1\end{psmallmatrix}\)，取
\[
\widetilde q=C^{-1}q,\qquad
\widetilde B=C^{-1}BC
=\begin{pmatrix}
a-uc&au+b-ucu-ud\\
c&cu+d
\end{pmatrix}.
\]
则 \((pB)C=(pC)\widetilde B\)，所以右作用被正确搬到新坐标。两种配对分别满足
\[
\widetilde p\,\widetilde q=pq,\qquad
\widetilde q\,\widetilde p=C^{-1}(qp)C.
\]
前者在 \(R\) 中保持原值，后者在 \(S\) 中随内自同构改变坐标。与外侧作用相容后，这正是\cref{b2:prop:matrix-row-column-bimodules}的两个张量同构的坐标变换。若 \(R\) 非交换，\(au\)、\(ucu\)、\(ud\) 等因子的次序均不可交换。
\end{solution}

\begin{exercise}\label{b2:ex:16-matrix-center}
设 \(R\) 为含幺环、\(n\ge1\)。直接用矩阵单位证明 \(Z\bigl(M_n(R)\bigr)=\bigl\{\,cI_n:c\in Z(R)\,\bigr\}\)。把这个显式同构与 Morita 中心不变量比较。
\end{exercise}
\begin{solution}
与各 \(E_{ii}\) 交换使中心矩阵的非对角条目全为零；再与 \(E_{ij}\) 交换，使全部对角元相同，记为 \(c\)。与在某个对角位置放任意 \(r\in R\) 的矩阵交换，得到 \(cr=rc\)，所以 \(c\in Z(R)\)。反向，中心元素乘单位矩阵与所有矩阵交换。因此 \(c\mapsto cI_n\) 是中心之间的环同构，正体现矩阵环与原环虽可能维数不同，中心仍相同。
\end{solution}

\begin{exercise}\label{b2:ex:16-center-obstruction}
设 \(k\) 为域。证明 \(M_2\bigl(k[\varepsilon]/(\varepsilon^2)\bigr)\) 与 \(M_2(k\times k)\) 不 Morita 等价，尽管它们作为 \(k\) 向量空间的维数相同。
\end{exercise}
\begin{solution}
两个环的 \(k\) 维数都为八。其中心分别为 \(k[\varepsilon]/(\varepsilon^2)\) 与 \(k\times k\)。前者有非零幂零元 \(\varepsilon\)，后者任何幂零元素的两个域坐标都必须为零，所以没有非零幂零元。因此中心不同构，由\cref{b2:prop:morita-center}，两个环不 Morita 等价。单独比较代数维数不能判断等价。
\end{solution}

\begin{exercise}\label{b2:ex:16-dimension-length-generators}
设 \(k\) 为域。在 \(k\) 与 \(S=M_3(k)\) 的标准等价中，求 \(V=k^7\) 的像的 \(k\) 维数、\(S\) 长度与最少 \(S\) 生成元数。该像是否为自由 \(S\) 模？
\end{exercise}
\begin{solution}
像为 \(V^3\cong(k^3)^7\)，其 \(k\) 维数为二十一。自然列模 \(k^3\) 简单，所以 \(S\) 长度为七，与 \(V\) 的 \(k\) 模长度相同。在\cref{b2:prop:semisimple-top-generators}中代入矩阵块大小三及简单模重数七，最少生成元数为 \(\lceil7/3\rceil=3\)，而原 \(k\) 模需要七个生成元。

它不是自由 \(S\) 模。有限秩自由 \(S\) 模的 \(k\) 维数为九的整数倍，二十一不是；无限秩自由模的 \(k\) 维数无限，也不可能。它仍然投射，因为半单环上的每个模投射，或因为等价保持原自由 \(k\) 模的投射性。
\end{solution}

\begin{exercise}\label{b2:ex:16-indecomposable-transport}
设 \(k\) 为域，令 \(R=k[t]\)、\(A=M_2(R)\)，把 \(M=R/(t^2)\oplus R/(t^3)\) 通过标准等价送到 \(A\) 模。求所得模的不可分解分量个数、长度及 \(k\) 维数，并解释为何不能从底 \(R\) 模的直和个数判断 \(A\) 不可分解性。
\end{exercise}
\begin{solution}
由\cref{b2:thm:module-correspondence}，\(R/(t^j)\) 的子模对应包含 \((t^j)\) 的理想，恰为 \((t^\ell)/(t^j)\)、\(0\le\ell\le j\)。这些子模构成一条链，相邻商均为 \(k\)，所以原循环模不可分解，长度分别为二和三。\cref{b2:prop:matrix-morita-equivalence}给出所用等价；全忠实性和\cref{b2:lem:module-equivalence-structure}的加性给出自同态环同构。由\cref{b2:prop:idempotent-direct-sum}，非平凡幂等元的不存在保证两者的像仍不可分解。再由\cref{b2:prop:morita-invariants}保持长度，像为两个不可分解 \(A\) 模的直和，长度为五。

底 \(k\) 维数乘二，故为十。每个像在忘掉矩阵作用后形如 \(\bigl(R/(t^j)\bigr)^2\)，但两个坐标会被矩阵单位相互送入；单个坐标不是 \(A\) 子模，不能将其当作 \(A\) 模的两个直和因子。
\end{solution}

\begin{exercise}\label{b2:ex:16-ring-chain-invariants}
设 \(k\) 为域、\(R=k[t]/(t^4)\)、\(A=M_3(R)\)。证明两环都是左 Noether 环和左 Artin 环，计算两个正则左模的合成长度。说明这是否与 Morita 等价保持模长度矛盾。
\end{exercise}
\begin{solution}
\(R\) 的理想为 \((t^j)\)、\(0\le j\le4\)，相邻商为 \(k\)，所以 \(\ell_R(R)=4\)。它满足两种链条件，由\cref{b2:prop:morita-invariants}及矩阵环等价，\(A\) 也满足两种链条件。

令 \(J=M_3((t))\)。\(J^j=M_3((t^j))\)，故
\[
A\supset J\supset J^2\supset J^3\supset0
\]
的四个层商各为 \(M_3(k)\) 的正则左模，分别是三个简单列模的直和。长度可加给出 \(\ell_A(A)=12\)。这没有矛盾：\cref{b2:prop:matrix-morita-equivalence}把 \(A\) 送到 \(R^3\)，其 \(R\) 长度正为十二；它把 \(R\) 送到自然列模 \(R^3\)，该模的 \(A\) 长度为四。等价保持对应对象的长度，两侧正则模却不是对应对象。
\end{solution}

\begin{exercise}\label{b2:ex:16-finite-presentation}
在 \(\mathbb Z\) 与 \(S=M_2(\mathbb Z)\) 的标准等价中，取 \(M=\mathbb Z/2\mathbb Z\oplus\mathbb Z/3\mathbb Z\)。将其像识别为 \((\mathbb Z/6\mathbb Z)^2\)，并写出一个只用有限自由 \(S\) 模的展示，明确检查全部关系。
\end{exercise}
\begin{solution}
中国剩余同构把 \((\bar a,\bar b)\) 送到 \(3a+4b\pmod6\)，逐坐标使用后得到题述识别。令 \(e=E_{11}\)，定义
\[
q:S\longrightarrow(\mathbb Z/6\mathbb Z)^2,\qquad
q(B)=\text{\(B\) 的第一列模六}.
\]
这是满的左 \(S\) 同态。其核中的第一列每项被六整除，第二列任意，因此
\[
\ker q=S(6e)+S(1-e).
\]
所以完整的有限自由展示为
\[
S^2\xrightarrow{\ d\ }S\xrightarrow{\ q\ }
(\mathbb Z/6\mathbb Z)^2\longrightarrow0,\qquad
d(B,C)=6Be+C(1-e).
\]
这里等式 \(\operatorname{im}d=\ker q\) 检查了全部关系，而不只是列出若干已知关系。\cref{b2:prop:morita-invariants}保证有限展示性保持，本题给出原坐标中的具体实现。
\end{solution}

\begin{exercise}\label{b2:ex:16-compose-bimodule-equivalences}
设 \(R,S,T\) 为含幺环，幺性双模 \({}_RP_S\)、\({}_SQ_T\) 分别通过张量实现从全部幺左 \(S\) 模到全部幺左 \(R\) 模、从全部幺左 \(T\) 模到全部幺左 \(S\) 模的等价。证明复合等价由 \({}_R(P\otimes_SQ)_T\) 实现，并说明该左 \(R\) 模为何为有限生成投射生成元。
\end{exercise}
\begin{solution}
对左 \(T\) 模 \(N\)，\cref{b2:thm:tensor-associativity-unit}给出自然同构
\[
P\otimes_S(Q\otimes_TN)\cong(P\otimes_SQ)\otimes_TN,
\qquad p\otimes(q\otimes n)\longmapsto(p\otimes q)\otimes n.
\]
因此复合确实由题述双模张量实现，作用次序不能颠倒。复合仍为范畴等价，而它将正则模 \(T\) 送到 \((P\otimes_SQ)\otimes_TT\cong P\otimes_SQ\)。由\cref{b2:thm:morita-converse}，正则模在等价下的这个像必为有限生成投射生成元，并有自同态反环同构于 \(T\)。
\end{solution}

\begin{exercise}\label{b2:ex:16-equivalent-to-division-ring}
设 \(D,E\) 为除环，\(m,n\ge1\)。证明 \(M_n(D)\) 与 \(M_m(E)\) Morita 等价，当且仅当 \(D\cong E\)。说明如何从唯一简单模的自同态环恢复除环，而不靠比较矩阵大小。
\end{exercise}
\begin{solution}
每个矩阵环只有一种简单左模类型，分别为自然列模 \(D^n\)、\(E^m\)。任意模范畴等价保持简单对象及其自同态环，所以给出
\[
\operatorname{End}_{M_n(D)}(D^n)
\cong\operatorname{End}_{M_m(E)}(E^m).
\]
与矩阵单位交换迫使自同态逐坐标相同；与任意 \(dE_{11}\) 交换，又使这一坐标映射为右乘某个除环元素。令 \(r_a\) 表示逐坐标右乘 \(a\)，则 \(r_a\circ r_b=r_{ba}\)，故这两个自同态环分别为 \(D^{\mathrm{op}}\)、\(E^{\mathrm{op}}\)。取反环得到 \(D\cong E\)。反向由矩阵环与原环的等价及改写标量得到。

因此\cref{b2:cor:morita-division-ring}中的除环同构类型可由简单对象恢复，而矩阵大小并不能由普通模范畴恢复。
\end{solution}

\begin{exercise}\label{b2:ex:16-full-corner-ideals}
设 \(k\) 为域、\(B=k[t]\)、\(R=M_3(B)\)、\(e=E_{11}+E_{22}\)。利用满角环的理想对应，求 \(R\) 的全部双边理想；计算 \(I=M_3((t^2))\) 对应的角理想及两个商环。
\end{exercise}
\begin{solution}
先对任意 \(n\ge1\) 考察 \(M_n(B)\) 的双边理想 \(L\)。若 \(X=(x_{ij})\in L\)，则
\[
E_{ai}XE_{jb}=x_{ij}E_{ab}\in L.
\]
因而各条目构成同一个 \(B\) 理想 \(\mathfrak a\)，且 \(L=M_n(\mathfrak a)\)；反向这类矩阵确为双边理想。\(B\) 是主理想整环，所以 \(R\) 的全部双边理想为 \(M_3((g))\)，其中 \(g\in k[t]\)，取零或首一多项式可消除单位倍的不唯一性。

这里 \(1=e+E_{31}eE_{13}\)，故 \(e\) 满，角环为 \(M_2(B)\)。\cref{b2:prop:full-corner-ideal-correspondence}把 \(M_3((g))\) 送到 \(M_2((g))\)，反向由 \(RJR\) 恢复。对题设 \(I\)，两个商环为
\[
R/I\cong M_3(k[t]/(t^2)),\qquad
eRe/eIe\cong M_2(k[t]/(t^2)).
\]
它们的矩阵大小不同，但满角在取商后仍满，故两商环 Morita 等价。
\end{solution}

\begin{exercise}\label{b2:ex:16-symmetric-alternating-forgotten}
在 \(V=\mathbb R^2\) 上，分别取矩阵为 \(I_2\) 的对称型与矩阵为 \(J=\begin{psmallmatrix}0&1\\-1&0\end{psmallmatrix}\) 的交替型。求两种伴随运算在 \(M_2(\mathbb R)\) 中的不动子空间维数，并说明为什么普通模范畴没有记录这一区别。
\end{exercise}
\begin{solution}
第一种不动条件是 \(X^{\mathsf T}=X\)，对称矩阵由两个对角条目和一个非对角条目确定，维数为三。第二种条件为 \(J^{-1}X^{\mathsf T}J=X\)。若 \(X=\begin{psmallmatrix}a&b\\c&d\end{psmallmatrix}\)，左边为 \(\begin{psmallmatrix}d&-b\\-c&a\end{psmallmatrix}\)，故不动条件迫使 \(a=d,b=c=0\)，只有标量矩阵，维数为一。底层实模和自同态环没有改变，改变的是另行指定的反转乘法运算；普通模范畴没有把它作为结构的一部分。
\end{solution}

\begin{exercise}\label{b2:ex:16-ideal-autoequivalence}
沿用\cref{b2:ex:krull-schmidt-nonunique}，令 \(s=\sqrt{-5}\)、\(R=\mathbb Z[s]\)、\(I=(2,1+s)\)。证明 \(I\) 是生成元且 \(\operatorname{End}_R(I)\cong R\)，从而给出 \(R\) 模范畴的一个 Morita 自等价。不得因 \(I\) 非自由便认为它不能作为生成元。
\end{exercise}
\begin{solution}
\cref{b2:ex:krull-schmidt-nonunique}已证明 \(I^2=2R\) 及 \(I\oplus I\cong R^2\)，所以由\cref{b2:prop:finite-projective-summand}，\(I\) 有限生成且投射。两个 \(R\) 线性映射 \(f_1(x)=2x\)、\(f_2(x)=(s-1)x/2\) 都从 \(I\) 取值于 \(R\)：第二个的分子属于 \(I^2=2R\)。并且
\[
f_1(2)+f_2(1+s)=4+\frac{s^2-1}{2}=1.
\]
所以 \((x,y)\mapsto f_1(x)+f_2(y)\) 给出 \(I\oplus I\) 到 \(R\) 的满同态，其截面为 \(r\mapsto\bigl(2r,(1+s)r\bigr)\)。因此 \(R\) 是 \(I\oplus I\) 的直和因子，\cref{b2:thm:generator-characterizations}说明 \(I\) 为生成元。

对任意 \(R\) 线性 \(f:I\rightarrow I\)，等式 \(2f(x)=f(2x)=x\cdot f(2)\) 说明它在分式域中为乘 \(c=f(2)/2\)。因 \(cI\subseteq I\)，乘以 \(I\) 得 \(cI^2\subseteq I^2\)，即 \(2cR\subseteq2R\)，消去非零整数二得 \(c\in R\)。反向每个 \(c\in R\) 都给出自同态，且这个对应保持加法、乘法与单位，所以 \(\operatorname{End}_R(I)\cong R\)。

\(R\) 交换，反环仍为 \(R\)。因为 \(I\) 是有限生成投射生成元，由\cref{b2:thm:morita-equivalence}，\(\operatorname{Hom}_R(I,-)\) 与 \(I\otimes_R-\) 通过上述环识别互为等价。这个生成元非自由，正说明自等价可以由非自由投射生成元实现。
\end{solution}
